\documentclass[11pt]{article}
\usepackage[utf8]{inputenc}
\usepackage{amsmath,amssymb,amsthm}
\usepackage{hyperref}
\usepackage{graphicx}
\usepackage{booktabs}
\usepackage{listings}
\usepackage{array}
\usepackage[margin=1in]{geometry}

\lstset{basicstyle=\ttfamily\small, breaklines=true, frame=single}

%% Bundle target: I1 (tier 3, journal CPC | Phys. Rep.)
%% Schema: docs/BUNDLE_DIRECTORY_SCHEMA.md
%% Source manifest: papers/I1/source_manifest.md
%% Append log: papers/I1/append_log.json
%% Change log: papers/I1/change_log.md
%% This bundle is a synthesis-driven new composition (Phase 7a sub-wave
%% 7a.2). It draws on paper15_methodology, paper44_riemannian_connection,
%% Phase 6f W1-W6 algebraic-GR substrate rosters, and Phase 6g W5
%% structural-Prop scoping pattern as substrate; the prose is authored
%% fresh.
% Auto-generated by update_counts.py — 2026-04-24T19:44:03.140400
% Do not edit manually. Run: uv run python scripts/update_counts.py
\newcommand{\totaltheorems}{3993}
\newcommand{\substantivetheorems}{3882}
\newcommand{\placeholdertheorems}{111}
\newcommand{\axiomcount}{1}
\newcommand{\sorrycount}{0}
\newcommand{\sorrypercent}{0.0}
\newcommand{\leanmodules}{167}
\newcommand{\leandefinitions}{3297}
\newcommand{\pythonmodules}{68}
\newcommand{\testfiles}{59}
\newcommand{\totaltests}{2836}
\newcommand{\figurecount}{111}
\newcommand{\notebookcount}{54}
\newcommand{\papercount}{19}
\newcommand{\aristotleproved}{322}
\newcommand{\aristotleruns}{44}
% --- Paper 18 / Phase 5t per-section counts (FermiHubbardDimer.lean) ---
\newcommand{\fhdTotal}{143}
\newcommand{\fhdWaveTwo}{13}
\newcommand{\fhdWaveThree}{6}
\newcommand{\fhdWaveFour}{22}
\newcommand{\fhdWaveFive}{22}
\newcommand{\fhdWaveSixABC}{38}
\newcommand{\fhdWaveSixExtra}{15}
\newcommand{\fhdWaveSeven}{19}
\newcommand{\fhdWaveEight}{8}

%% ── Per-module Phase-6f/6g count-macro fallbacks ────────────────────
%% These 19 macros are referenced in §11–§13 but are not (yet) emitted
%% by scripts/update_counts.py into docs/counts.tex. Values computed
%% from lean/lean_deps.json (kind == "theorem" grouped by module — the
%% exact update_counts.py convention), snapshot 2026-06-10.
%% \providecommand yields to docs/counts.tex definitions the moment the
%% counts pipeline emits these macros; the durable fix is to add them
%% to update_counts.py and delete this block.
\providecommand{\curvatureThms}{14}
\providecommand{\einsteinTensorThms}{9}
\providecommand{\energyConditionsThms}{8}
\providecommand{\exactSolutionsThms}{16}
\providecommand{\admFormalismThms}{10}
\providecommand{\tetradFormalismThms}{5}
\providecommand{\phaseSixfWonesixTotal}{62}
\providecommand{\phaseSixfWsevenTotal}{13}
\providecommand{\riemannCoordinateThms}{8}
\providecommand{\riemannDifferentialBianchiThms}{8}
\providecommand{\bundleRiemannAuxThms}{6}
\providecommand{\bundleRiemannThms}{7}
\providecommand{\leviCivitaThms}{7}
\providecommand{\phaseSixfWeightTotal}{36}
\providecommand{\riccatiComparisonThms}{11}
\providecommand{\mazurSigmaModelRigidityThms}{12}
\providecommand{\hawkingPenroseCurveThms}{6}
\providecommand{\areaTheoremCurveThms}{10}
\providecommand{\phaseSixgWnineTotal}{39}
%% ── end count-macro fallbacks ───────────────────────────────────────

\title{Continuous Verification for Computational Physics Research:\\
A Three-Layer Pipeline with Worked Cases from Schwinger--Keldysh
Effective Field Theory, Algebraic General Relativity, and
Topological Quantum Field Theory}

\author{John Roehm}
\date{Draft: Phase 7a, May 2026}

\begin{document}

\maketitle

\begin{abstract}
Computational physics research lacks a structural guarantee that
published numerical claims match the underlying computation. We
describe a three-layer continuous-verification pipeline integrating
Python numerics, Lean~4 formal proofs, and the Aristotle automated
theorem prover, organized into a fourteen-stage workflow with
fifteen invariants. Applied across a research program spanning
dissipative effective field theory, analog Hawking radiation,
emergent gravity through black-hole thermodynamics, anomaly
constraints on Standard Model particle content, dark-sector substrate
constraints, and topological quantum computation foundations, the
pipeline has produced \totaltheorems{} machine-checked theorems
across \leanmodules{} Lean modules with \sorrycount{} remaining gaps
(\sorrypercent\%) plus \aristotleproved{} theorems closed by
automated theorem proving. We document three new pipeline stages
beyond the original twelve-stage version (interactive Lean
development via the language-server protocol, fresh-context
adversarial review, and meta-process quality improvement), three new
invariants (placeholder non-load-bearing, no heartbeat overrides in
proof bodies, primary-source citation cache), the publication-bundle
architecture that consolidates per-wave drafts into the
publication targets that frame enumerates, and the preemptive-strengthening discipline
that reduces theorem-statement audit cost from twelve retroactive
cuts to five and then to zero across three successive waves (with
a mandatory ruthless post-wave review as a safety net for residual
subtler patterns). Three worked cases illustrate the methodology:
the Aristotle FirstOrderKMS counterexample (where automated proving
failed and the failure surfaced a substantive physics result), the
gap-solution-bounded conjecture disproof (where formalization-driven
counterexample search refuted a literature claim), and the
chirality-wall axiom decomposition (where structural axiom
classification turned a single 50-term proof into nine
individually-tractable sublemmas). Two additional sidebar studies
document the algebraic-general-relativity substrate formalization
work that the pipeline accreted as a byproduct, and a
Mathlib-dependency-fallback discipline pattern for theorems whose
substantive content has no upstream PDE infrastructure available.
A final process result documents the limits of metadata-level
citation validation: a June~2026 external-review remediation found
that a citation can be metadata-valid (resolvable identifier,
cached primary source, correct bibliographic data) while the
specific claim the prose attributes to it is unsupported by the
source (three instances in roughly ten spot checks), motivating
three new deterministic prose-consistency checks and a systematic
content-level attribution sweep.
The methodology is domain-agnostic.
\end{abstract}

%% ─────────────────────────────────────────────────────────────────
%% §1 Introduction
%% ─────────────────────────────────────────────────────────────────

\section{Introduction}
\label{sec:intro}

Numerical computation in physics research is increasingly
sophisticated and increasingly opaque. A typical large
computational-physics program now combines several thousand lines of
Python or Julia, dozens of Jupyter notebooks, and a paper-trail of
PDF drafts whose numerical content is the output of code paths that
the published prose does not name. The standard publication workflow
provides no machine-checkable connection between the prose and the
code: reviewers cannot, and do not, run the computation. Five years
after publication, when the original code has bit-rotted and the
authors have rotated, no one can reconstruct which version of which
formula produced which figure.

This is not a hypothetical concern. We have personally observed
each of the following failure modes during a single research
program: (i) a numerical claim in a draft that traced to a code
path that no longer existed; (ii) a ``formally verified'' claim
whose underlying Lean theorem was a placeholder of type
\texttt{True}; (iii) a citation whose arXiv identifier resolved to
an unrelated graph-neural-network paper; (iv) a parameter whose
value was inherited from a prior draft and had drifted from the
primary source it cited. None of these were detected by the
authors before adversarial review. All of them shipped to draft
state.

The pipeline we describe makes verification \emph{continuous}, in
the sense of integrated into the daily research workflow rather
than retrofitted before submission. Antecedents on the
formal-verification side include landmark engineering verifications
(CompCert~\cite{compcert} for a C compiler; seL4~\cite{sel4} for an
operating-system kernel) and the more recent emergence of
Lean-based physics infrastructure
(PhysLean~\cite{physlean}). Antecedents on the
proof-assistance side include large-language-model copilots for
Lean theorem proving~\cite{leancipilot} and recent
silver-medal IMO results from automated proving~\cite{alphaproof}.
Our pipeline has fourteen stages and fifteen invariants, integrates
three orthogonal verification layers, and produces
machine-checkable artifacts at each stage. We have applied it to
a research program spanning \papercount{}
papers covering Schwinger--Keldysh dissipative effective field
theory, analog Hawking radiation in Bose--Einstein condensates and
graphene, emergent gravity through black-hole thermodynamics,
anomaly constraints on Standard Model particle content, dark-sector
substrate constraints from causal-set theory, entropic gravity, and
Jacobson thermodynamic emergence, and topological quantum
computation foundations. The pipeline has produced
\totaltheorems{} machine-checked theorems across \leanmodules{}
Lean modules, with \sorrycount{} remaining unproved gaps
(\sorrypercent\%) and \aristotleproved{} theorems closed by the
Aristotle automated theorem prover. The pipeline, its process
documents, and the validation machinery described below are
publicly available; monospaced file paths throughout this paper
are paths from the root of the public repository (see the Code
Availability statement).

This paper extends and supersedes an earlier internal draft
description of a twelve-stage version of the pipeline (April
2026). The version described here adds three stages beyond
the original twelve: interactive Lean development via the
language-server protocol (Stage~3a), fresh-context adversarial
review (Stage~13), and meta-process quality improvement
(Stage~14); and three invariants (placeholder non-load-bearing,
no heartbeat overrides in proof bodies, primary-source citation
cache). The invariant set has since grown to fifteen; the
fifteenth (explicit user sign-off, with a documented discharge
plan, for any new project-local axiom) encodes the project's
posture that axioms are temporary scaffolding, not permanent
commitments (Section~\ref{sec:pipeline}). It also documents the
publication-bundle architecture that consolidates per-wave drafts
into the publication targets that frame enumerates, and the preemptive-strengthening
discipline that reduces theorem-statement audit cost from twelve
retroactive cuts (Phase~6c.3) to five (Phase~6b.1) and then to zero
(Phase~6m R6) by applying strengthening rules prospectively, with
a mandatory ruthless post-wave review as the safety net for the
subtler patterns the prospective checklist misses.

The remainder of this paper is organized as follows.
Section~\ref{sec:architecture} describes the three-layer
verification architecture (Python numerics, Lean~4 formal proofs,
Aristotle automated proving). Sections~\ref{sec:case1},
\ref{sec:case2}, \ref{sec:case3} present three worked cases that
illustrate the pipeline catching errors that single-layer
verification cannot. Section~\ref{sec:pipeline} presents the
fourteen-stage workflow and its invariants in detail.
Section~\ref{sec:strengthening} documents the preemptive-strengthening
discipline.
Section~\ref{sec:provenance} describes sentence-level provenance and
the cross-paper consistency cluster mechanism that enables bundle
publication. Section~\ref{sec:three_layer_extended} returns to the
three-layer architecture in the context of automated theorem proving
integration. Section~\ref{sec:adversarial} describes the
fresh-context adversarial-reviewer pattern, the recorded-review
readiness semantics, and the limits of metadata-level citation
validation surfaced by a June~2026 external-review remediation.
Section~\ref{sec:phase6f_substrate} presents the
algebraic-general-relativity substrate roster shipped over
Phase~6f Waves~1--6. Section~\ref{sec:riemannian_sidebar}
presents the Phase~6f Wave~7 catch-up Riemannian connection sidebar
and its bundle-level extension via the abstract-bundle Levi-Civita
construction (Phase~6f Wave~8 Sessions~1--3).
Section~\ref{sec:structural_prop_scoping} describes the
Mathlib-dependency-fallback structural-Prop scoping discipline
(Phase~6g Wave~5). Section~\ref{sec:lessons} discusses lessons
learned and outlines future pipeline extensions.

%% ─────────────────────────────────────────────────────────────────
%% §2 Three-layer verification architecture
%% ─────────────────────────────────────────────────────────────────

\section{The Three-Layer Verification Architecture}
\label{sec:architecture}

The pipeline integrates three orthogonal verification layers, each
addressing a failure mode that the others cannot.

\textbf{Layer~1: Python numerics for fast iteration.} A formula is
implemented in a canonical \texttt{formulas.py} module, with a
docstring containing (i) the mathematical formula in plain text,
(ii) the Lean theorem name that formalizes it, (iii) the Aristotle
run identifier (or ``manual'' or ``pending''), and (iv) a citation
to the published source for derived formulas. Domain-specific
modules import from \texttt{formulas.py} rather than reimplementing.
Python tests verify computational correctness, edge cases, and
physical bounds. Layer~1 catches computation errors but not
formalization errors.

\textbf{Layer~2: Lean~4 formal proofs for mathematical certainty.}
Each formula is stated as a Lean theorem in the strongest form
possible. Statements are not weakened for proof convenience; if
the strongest form is not provable, the proof architecture is
wrong (Section~\ref{sec:strengthening}). Manual proofs are written
for theorems within reach. Layer~2 catches formalization errors
but not definition errors: a formally verified proof can still
be wrong if the definitions it operates on are wrong
(Section~\ref{sec:case3}).

\textbf{Layer~3: Aristotle automated theorem prover for sorry
gaps.} The entire Lean project is submitted with strategy hints
embedded in the docstrings of unproved theorems (the ``PROVIDED
SOLUTION'' convention). Aristotle uses tactic synthesis, term
search, and machine learning to find proofs. Results are
manually reviewed before integration (no blind acceptance)
to verify that the proof exercises the hypotheses, that
\texttt{lake build} passes after integration, and that no
heartbeat overrides have been introduced (Invariant~10).

The three layers interact in a discovery--formalization--automation
cycle. A typical cycle: a Python implementation of a new formula
exposes a numerical regularity; the regularity is stated as a Lean
theorem with sorry; the theorem is closed by Aristotle in a
priority batch; the closure feeds back into Python by setting the
Aristotle run identifier in the formula docstring, which is then
checked by the validation suite (Section~\ref{sec:pipeline}). The
critical-coupling theorem $G_c = 8\pi^2/(N_f \Lambda^2)$ for the
chirality-wall mechanism completed all three layers in a single
session: derived in Python, stated in Lean, closed by Aristotle in
run \texttt{79e07d55}, and checked by validation, all within
twenty-four hours.

The composition is asymmetric in a way that matters. Layer~1 alone
gives reproducibility but not correctness: notebooks can produce
the same wrong answer indefinitely. Layer~2 alone gives correctness
but not numerical scope: a formalization of an idealized model is
silent on the actual experimental parameters.
Layer~3 alone gives proof-search throughput but no verification of
what was proved: an automated proof can be vacuous if the
definitions it operates on are wrong (Section~\ref{sec:case3}). The
composition gives all three.

\begin{figure}[h]
\centering
\includegraphics[width=0.85\linewidth]{figures/i1_fig1_three_layer_architecture.png}
\caption{Three-layer verification architecture (this paper, \S2).
Layer~1 (Python numerics, top, blue) feeds Layer~2 (Lean~4 formal
proofs, middle, amber) feeds Layer~3 (Aristotle automated theorem
prover, bottom, sage). The curved feedback arrow on the right
represents counterexample- or refinement-driven flow back to Layer~1
(e.g.\ Aristotle run~\texttt{270e77a0} returned the counterexample
that forced the strengthened algebraic \texttt{FirstOrderKMS}
condition; see Section~\ref{sec:case1}). The
asymmetric composition is the substantive structural claim of \S2:
none of the three layers alone suffices.}
\label{fig:three_layer}
\end{figure}

%% ─────────────────────────────────────────────────────────────────
%% §3 Worked case 1: FirstOrderKMS Aristotle counterexample
%% ─────────────────────────────────────────────────────────────────

\section{Worked Case~1: The FirstOrderKMS Aristotle Counterexample}
\label{sec:case1}

An automated prover can return a refutation instead of a proof, and the
refutation is the more valuable of the two outcomes. This section works
through the case in which that happened here. A uniqueness theorem for
the first-order Schwinger--Keldysh action was submitted with a
hypothesis that had been written to encode the Kubo--Martin--Schwinger
(KMS) condition and did not encode it. Aristotle returned an explicit
nine-component coefficient vector satisfying the hypothesis and
violating the conclusion. The defect was in the hypothesis, not in the
physics being formalized and not in the prover: the informal sentence
the hypothesis was meant to capture (``KMS fixes the first-order
dissipative action up to two transport coefficients'') is true, and is
now the closed theorem
\texttt{SKEFTHawking.SKDoubling.firstOrder\_uniqueness}
(\texttt{lean/SKEFTHawking/SKDoubling.lean}, lines~412--439, Aristotle
run~\texttt{270e77a0}). What was wrong was the formal rendering of one
of its three hypotheses, and nothing short of a machine attempting the
proof was going to say so.

\subsection{The formal setting}

Two distinct nine-fold objects appear in this case, and keeping them
apart is the whole of the bookkeeping. The first is the local field
content. \texttt{SKFields} (\texttt{SKDoubling.lean}, line~80) is a
record of nine real numbers: the retarded and advanced fields
$\psi_r, \psi_a$; their first derivatives
$\partial_t\psi_r, \partial_x\psi_r, \partial_t\psi_a, \partial_x\psi_a$;
and the second derivatives
$\partial_t^2\psi_r, \partial_t\partial_x\psi_r, \partial_x^2\psi_r$.
An \texttt{SKAction} is a map from this record to a pair
$(\operatorname{Re}\mathcal{L}, \operatorname{Im}\mathcal{L})$ of reals,
so hermiticity of the Lagrangian is carried by the data type rather
than by an assumption.

The second nine-fold object is the coefficient space.
\texttt{FirstOrderCoeffs} (lines~292--301) enumerates nine candidate
monomials, six in the retarded-response sector and three in the noise
sector:
%
\begin{align*}
\operatorname{Re}\mathcal{L} &=
  r_1\,\psi_a\partial_t^2\psi_r
+ r_2\,\psi_a\partial_x^2\psi_r
+ r_3\,\psi_a\partial_t\partial_x\psi_r
+ r_4\,\psi_a\partial_t\psi_r
+ r_5\,\psi_a\partial_x\psi_r
+ r_6\,\psi_a\psi_r , \\
\operatorname{Im}\mathcal{L} &=
  i_1\,\psi_a^2
+ i_2\,(\partial_t\psi_a)^2
+ i_3\,(\partial_x\psi_a)^2 ,
\end{align*}
%
assembled by \texttt{firstOrderAction} (lines~304--315). The target of
the uniqueness statement is the two-parameter family
\texttt{firstOrderDissipativeAction} (lines~251--261), built from
\texttt{DissipativeCoeffs} (\texttt{lean/SKEFTHawking/Basic.lean},
lines~152--158), which carries $\gamma_1$ (isotropic phonon damping,
related to bulk viscosity), $\gamma_2$ (anisotropic damping along the
superfluid flow), and the two non-negativity fields
\texttt{gamma\_1\_nonneg} and \texttt{gamma\_2\_nonneg}. Uniqueness
means: every point of the nine-dimensional coefficient space that
satisfies the constraints is the image of some $(\gamma_1,\gamma_2)$.

\subsection{The published axiomatic and its Lean rendering}

The constraints are not this program's invention. Crossley, Glorioso,
and Liu~\cite{CrossleyGloriosoLiu2017} impose on the closed-time-path
action, in their \S I.D, a normalization condition (their eq.~(1.28),
sharpened at the level of the action to eq.~(1.36): the action vanishes
when the two contour legs are identified, so every term carries at
least one factor of the advanced field), a $\mathbb{Z}_2$ reflection
condition on the action (their eq.~(1.29)), and, as item~5 of the same
list, the requirement that the imaginary part of the action be
non-negative. For a thermal density matrix (their eq.~(1.60)) the
generating functional obeys in addition the local KMS condition, a
$\mathbb{Z}_2$ symmetry (their \S I.F, eq.~(1.68)) whose quadratic-order
content is the familiar fluctuation--dissipation theorem
$\operatorname{Im}G_R = \tanh(\beta_0\omega/2)\,G_S$ (their eq.~(1.62)).
Glorioso, Crossley, and Liu~\cite{GloriosoCrossleyLiu2017} replace the
local KMS condition by the \emph{dynamical} KMS symmetry (their \S IV.A,
eq.~(4.1)), equivalent to it in the classical limit but applicable to
any state in local equilibrium rather than only to a thermal density
matrix perturbed by background fields. Glorioso and
Liu~\cite{GloriosoLiu2018} then derive the local second law
$\partial_\mu S^\mu \ge 0$ from that $\mathbb{Z}_2$ symmetry together
with the classical remnant of quantum unitarity, which supplies both
$\operatorname{Im}I_{\mathrm{eff}} \ge 0$ and the normalization
condition (their \S II, eqs.~(2.5)--(2.6)); the second law is proved in
general and its local form perturbatively in the derivative expansion.
A pedagogical account of the same construction is given in the lectures
of Glorioso and Liu~\cite{GloriosoLiu2018TASI}. The realization of
dynamical KMS is not unique when the theory carries ultraviolet content
that the infrared symmetries do not fix: Jain and
Kovtun~\cite{JainKovtun2024} exhibit distinct Schwinger--Keldysh models
that share an effective action but realize the symmetry differently
(their Appendix~C, eqs.~(C.9)--(C.11)), which is why this program keeps
the strict and algebraic readings of the condition as separate
predicates rather than collapsing them.

Three of these conditions transcribe into Lean without difficulty.
Normalization is \texttt{satisfies\_normalization} (lines~126--129),
positivity of the imaginary part is \texttt{satisfies\_positivity}
(lines~140--141), and hermiticity is discharged by the
$\mathbb{R}\times\mathbb{R}$ codomain. The fourth condition is where the
case begins.

\subsection{The hypothesis as first submitted}

The KMS condition was first rendered at the level of the field
transform, as the structure \texttt{KMSSymmetry} (lines~156--167). It
bundles a map on \texttt{SKFields}, a specification of that map, and an
invariance clause. The specification is a conjunction of exactly four
equations, fixing $\psi_r$, $\partial_t\psi_r$ and $\partial_x\psi_r$
and shifting $\psi_a \mapsto \psi_a + \beta\,\partial_t\psi_r$. The
remaining five components of the record, namely
$\partial_t\psi_a$, $\partial_x\psi_a$ and the three second derivatives
of $\psi_r$, are named nowhere in the specification, so a witness may
transform them however it likes. That the specification pins exactly
four components is not a reading of the source but a theorem,
\texttt{GloriosoLiu.kmsSpec\_pins\_four\_components}
(\texttt{lean/SKEFTHawking/GloriosoLiu/Phase1Reconciliation.lean},
lines~118--127), whose proof projects the four conjuncts out of the
structure field itself.

Set against the published condition, this rendering fails in two
independent directions at once. It is too permissive, in that it
constrains four of nine field components and is silent on the rest. It
is also too rigid: its \texttt{invariance} clause demands pointwise
equality of the Lagrangian under the transform, whereas the dynamical
KMS condition of ref.~\cite{GloriosoCrossleyLiu2017} is equality up to a
total derivative (their eq.~(5.3)), imposed on combinations that are
closed under the symmetry rather than on individually truncated orders
(their eq.~(5.9) and the sentence following it). The library records
this by keeping the transform-level reading as
\texttt{GloriosoLiu.hasDynamicalKMS\_strict}
(\texttt{lean/SKEFTHawking/GloriosoLiu/Axioms.lean}, lines~91--92), a
sanity primitive used for degenerate witnesses, while the substantive
\texttt{dynamical\_KMS} field of the six-condition skeleton
\texttt{SKEFTAxioms} (lines~160--172) is
\texttt{hasDynamicalKMS\_algebraic} (lines~108--111). Neither reading is
a Lean \texttt{axiom} declaration: both are ordinary structures, and the
theorems below are proved from them.

\subsection{What came back}

Submitted with the transform-level hypothesis, the uniqueness statement
was refuted. The witness is the coefficient vector
%
\[
c = (r_1,\ldots,r_6,i_1,i_2,i_3) = (0,0,0,0,0,0,0,1,0),
\]
%
a single unit of noise in the $(\partial_t\psi_a)^2$ channel and nothing
else, now recorded as
\texttt{GloriosoLiu.aristotleCounterexample} (lines~139--140). It
defeats the statement for a reason the reader can check against
\texttt{firstOrderAction} directly. Its action is
$(\operatorname{Re}\mathcal{L},\operatorname{Im}\mathcal{L}) =
(0,(\partial_t\psi_a)^2)$, so positivity holds. Matching the imaginary
part against \texttt{firstOrderDissipativeAction}, whose noise sector is
$(\gamma_1/\beta)\psi_a^2 + (\gamma_2/\beta)(\partial_t\psi_a)^2$, forces
$\gamma_1 = 0$ and $\gamma_2 = \beta > 0$; but then the retarded sector
of the target is $\beta\,\psi_a\partial_t^2\psi_r$, which does not
vanish identically, while the witness has no retarded part at all. No
pair $(\gamma_1,\gamma_2)$ reproduces it. The coefficient $i_2$ that the
witness excites multiplies $(\partial_t\psi_a)^2$, and
$\partial_t\psi_a$ is precisely one of the five components the transform
specification never mentions.

The witness survives in the library as a kernel-checked object rather
than as an anecdote. \texttt{GloriosoLiu.aristotle\_counterexample\_violates\_FirstOrderKMS}
(lines~150--158) proves
$\neg\,\texttt{FirstOrderKMS}\ c\ \beta$ for every $\beta > 0$, and
\texttt{GloriosoLiu.four\_of\_nine\_partition\_recovered}
(lines~189--204) combines that with the four-component projection into
one statement: any instance of the transform-level condition pins those
four components, and the witness is ruled out only by the algebraic
condition. A corollary,
\texttt{GloriosoLiu.four\_of\_nine\_partition\_under\_GLU}
(lines~231--239), derives the same conclusion from the six-condition
skeleton by extracting the coefficient witness from its
\texttt{dynamical\_KMS} field. We emphasize the scope: what is refuted
is a formalization, not the axiomatic of
refs.~\cite{CrossleyGloriosoLiu2017,GloriosoCrossleyLiu2017}.

\subsection{The corrected condition}

The repair is to impose the KMS condition on the coefficients rather
than on the transform. \texttt{FirstOrderKMS} (lines~367--379) is a
\texttt{Prop}-valued structure carrying seven relations: the four
lower-derivative retarded monomials vanish,
$r_3 = r_4 = r_5 = r_6 = 0$; two fluctuation--dissipation relations tie
the surviving noise coefficients to the retarded ones,
$i_1\beta = -r_2$ and $i_2\beta = r_1 + r_2$; and the spatial noise
channel is switched off, $i_3 = 0$. The two fluctuation--dissipation
relations are stated multiplicatively in $\beta$, which is equivalent to
the familiar $i_1 = -r_2/\beta$ and $i_2 = (r_1+r_2)/\beta$ at any
positive temperature and avoids a division in the hypothesis. Every one
of the nine coefficients appears in at least one of the seven relations,
and the solution set is two-dimensional, coordinatized by $(r_1,r_2)$.

Under this hypothesis the uniqueness theorem closes.
\texttt{firstOrder\_uniqueness} states that for every
\texttt{FirstOrderCoeffs} $c$ and every $\beta > 0$, if
\texttt{firstOrderAction}~$c$ satisfies positivity and $c$ satisfies
\texttt{FirstOrderKMS}~$c$~$\beta$, then there exist
\texttt{DissipativeCoeffs} whose action agrees with it on every field
configuration. The proof supplies the witness explicitly,
$\gamma_1 = -r_2$ and $\gamma_2 = r_1 + r_2$, and discharges the two
non-negativity obligations of \texttt{DissipativeCoeffs} from positivity
together with the two fluctuation--dissipation relations. Exactly one of
the seven relations is what excludes the counterexample: with
$r_1 = r_2 = 0$ and $i_2 = 1$, the relation $i_2\beta = r_1 + r_2$
collapses to $\beta = 0$.

\subsection{Probing the corrected condition on purpose}

A hypothesis that has just been caught being too weak invites the
question of whether it is now too strong, and the same instrument
answers it. Two further statements were submitted deliberately. The
first, \texttt{firstOrder\_KMS\_optimal} (lines~582--603), proves that
under \texttt{FirstOrderKMS} and $\beta > 0$ the positivity of
\texttt{firstOrderAction}~$c$ is \emph{equivalent} to
$i_1 \ge 0 \wedge i_2 \ge 0$, so positivity contributes exactly the two
sign conditions that \texttt{DissipativeCoeffs} already carries and no
further relation among the coefficients is hiding in it. The second
mis-states the condition on purpose. \texttt{FirstOrderKMS\_altSign}
(lines~610--618) flips one sign, replacing $i_1\beta = -r_2$ by
$i_1\beta = +r_2$, and
\texttt{firstOrder\_altSign\_uniqueness\_test} (lines~636--656) is the
kernel-checked \emph{negation} of the corresponding uniqueness claim,
witnessed by $c = (1,1,0,0,0,0,1,2,0)$ at $\beta = 1$: that vector
satisfies the flipped relations and positivity, yet forces
$\gamma_1 = -r_2 = -1$, contradicting \texttt{gamma\_1\_nonneg}. The
sign in the fluctuation--dissipation relation is therefore load-bearing
rather than conventional, and the library says so in a form a reader can
recompute rather than in a remark.

\begin{figure}[h]
\centering
\includegraphics[width=0.95\linewidth]{figures/i1_fig4_firstorderkms_grid.png}
\caption{The nine first-order coefficients of \texttt{FirstOrderCoeffs}
under the two candidate KMS hypotheses. Left panel: under the
transform-level \texttt{KMSSymmetry}, whose specification pins four of
the nine \texttt{SKFields} components and leaves $\partial_t\psi_a$
among the five it does not mention, the noise coefficient $i_2$ carries
no algebraic constraint; the cell highlighted in carmine is the refuting
witness $c = (0,0,0,0,0,0,0,1,0)$ (\texttt{aristotleCounterexample}),
which satisfies positivity and is the action of no
$(\gamma_1,\gamma_2)$ pair. Right panel: the algebraic
\texttt{FirstOrderKMS} imposes seven relations, every coefficient
entering at least one of them and $i_3 = 0$ appearing explicitly,
leaving a two-dimensional solution set coordinatized by $(r_1,r_2)$ and
equivalently by $(\gamma_1,\gamma_2)$;
\texttt{firstOrder\_uniqueness} closes under it (Aristotle
run~\texttt{270e77a0}).}
\label{fig:firstorderkms_grid}
\end{figure}

\subsection{What generalizes}

The case supports a narrow claim and not a broad one. A prover's failure
to close a statement is evidence about the statement, and this
particular failure localized the defect to a single hypothesis out of
three and, within that hypothesis, to the five field components its
specification never named. That is more information than a human review
of the same statement produced, because the informal sentence the
statement encodes is true; the error lived entirely in the gap between
the sentence and its transcription, which is the one place prose review
cannot look. Three conditions made the outcome usable. The statement was
posed over a finite-dimensional coefficient space, so a counterexample
was expressible as a tuple rather than as an obstruction argument. The
refuting witness was retained and formalized, so the negative result is
now checkable by the kernel instead of remembered. And the corrected
hypothesis was afterwards attacked on purpose, by flipping a sign and
proving the resulting claim false, so its non-triviality is measured
rather than assumed.

The limits are equally concrete. The refutation did not supply
\texttt{FirstOrderKMS}; it established only that the submitted
hypothesis admitted a vector the conclusion could not accommodate.
Choosing the seven algebraic relations required reading the derivation
in refs.~\cite{CrossleyGloriosoLiu2017,GloriosoCrossleyLiu2017} and
projecting it onto this coefficient space by hand. A prover localizes a
defect; it does not supply the physics that repairs it. Nor does a
refutation distinguish between the two ways a statement can be wrong. A
hypothesis can be too weak, as here, or the definitions underneath a
correct-looking statement can themselves be wrong, in which case a
successful proof is the dangerous outcome rather than the reassuring one
(Section~\ref{sec:case3}). Only the first failure mode announces itself
by refusing to close.

\section{Worked Case~2: The Gap-Solution-Bounded Conjecture Disproof}
\label{sec:case2}

Worked Case~1 turned on a prover's failure to close a theorem that was
true only under a stronger hypothesis. This case is the sharper
variant. The target statement was false, the prover said so with an
explicit witness, and the statement that replaced it is weaker in a way
that changes what the surrounding physics is entitled to assume. The
episode is the one that most directly answers the objection that formal
verification only re-certifies what the author already believes:
nothing in the informal derivation was going to produce this
counterexample, because the informal derivation never wrote down the
hypothesis the counterexample violates.

% TODO: the expectation Delta < Lambda is stated below without
% attribution to any published work. It is presented as the
% EFT-validity expectation that the formalization plan encoded as a
% target theorem, which is documented. No primary source read for this
% section asserts the bound, so no citation is attached; an earlier
% framing that credited it to "the published literature on NJL gap
% equations" is not supported by any source obtainable from the
% citation cache. Retrieval would be needed to restore an attribution.

The emergent-gravity sector of the program formalizes a
Nambu--Jona-Lasinio-type self-consistency condition for tetrad
condensation, relating a condensate gap $\Delta$, a four-fermion
coupling $G$, a UV cutoff $\Lambda$, and a fermion flavour count $N_f$
(\texttt{lean/SKEFTHawking/TetradGapEquation.lean}). The gap equation
is the fixed-point condition
\begin{equation}
  \Delta \;=\; G\,N_f\,\Delta\,I(\Delta),
  \qquad
  I(\Delta) \;=\; \int_0^{\Lambda} \frac{c_4\,p^3}{p^2+\Delta^2}\,\mathrm{d}p
  \;=\; \frac{c_4}{2}\left[\Lambda^2 - \Delta^2
        \log\!\left(1+\frac{\Lambda^2}{\Delta^2}\right)\right],
  \label{eq:gapeq}
\end{equation}
with density-of-states coefficient $c_4 = 1/(4\pi^2)$ in $d=4$. The
right-hand evaluation is exact, and it is what the Lean development
takes as primitive: \texttt{gapIntegral} is defined as the closed form
rather than as a Bochner integral, which keeps the module inside real
analysis and away from measure theory. Two properties of $I$ carry the
whole case. It is strictly positive for every $\Delta \geq 0$ and every
$\Lambda > 0$ (\texttt{gapIntegral\_pos}), and it is strictly
decreasing in $\Delta$ (\texttt{gapIntegral\_strictAnti}). Neither
statement carries an upper restriction on $\Delta$, and neither
needs one: Eq.~(\ref{eq:gapeq}) is a convergent integral over a fixed
momentum range whatever the value of the gap.

The rest of the module is the standard phase-transition picture.
Writing $G_c = 1/(N_f\,I(0)) = 8\pi^2/(N_f\Lambda^2)$
(\texttt{criticalCoupling\_formula}), the trivial solution $\Delta = 0$
is the only solution below $G_c$
(\texttt{gap\_trivial\_unique\_subcritical}), a positive fixed point
exists above it (\texttt{gap\_nontrivial\_exists}, by the intermediate
value theorem applied to $g(\Delta) = G N_f I(\Delta) - 1$), and the
bifurcation is a pitchfork at $G = G_c$ (\texttt{bifurcation\_at\_Gc}).
Nine theorems of this module, including all of the above, closed in
Aristotle run~\texttt{79e07d55}.

A tenth was submitted with them. Above the critical coupling the gap is
a low-energy parameter of an effective description valid up to
$\Lambda$, so a gap exceeding the cutoff is outside the regime in which
the equation was derived, and the expectation that a solution obeys
$\Delta < \Lambda$ is the natural one. The formalization plan recorded
it as a target theorem, \texttt{gap\_solution\_bounded}, and stated it
the only way a Lean statement can be stated, with every hypothesis
enumerated: $0 < G$, $0 < N_f$, $0 < \Lambda$, $0 < \Delta$, and
$\Delta$ a fixed point of the gap operator, concluding $\Delta <
\Lambda$. That enumeration is the substantive difference from the
informal argument. On paper, ``the supercritical regime'' silently
carries a second condition, that $G$ is not too far above $G_c$,
because that is the regime anyone has in mind while writing. In Lean
the second condition either appears in the hypothesis list or it does
not exist, and here it did not.

Aristotle returned a witness rather than a proof: $N_f = 1$,
$\Lambda = 1$, $\Delta = 1$, and
$G = 1/\bigl[(c_4/2)(1-\log 2)\bigr]$. The check is two lines. For
$\Delta > 0$ the fixed-point condition of Eq.~(\ref{eq:gapeq}) reduces
to $G N_f I(\Delta) = 1$, and at $\Delta = \Lambda$ the closed form
gives $I(\Lambda) = (c_4\Lambda^2/2)(1-\log 2) = I(0)\,(1-\log 2)$, so
that value of $G$ makes $\Delta = \Lambda$ an exact solution. The
witness saturates the claimed strict inequality rather than violating
it by a margin, which is the shape a counterexample to a strict bound
usually takes and the shape a numerical scan is least likely to find.

The saturating witness is not a boundary artefact, and this is where
the formalization pays for itself a second time. Because $I$ is
strictly positive at every $\Delta$, the coupling
\begin{equation}
  G(\Delta) \;=\; \frac{1}{N_f\,I(\Delta)}
  \label{eq:invertgap}
\end{equation}
is well defined and positive for every $\Delta > 0$, and by
construction $\Delta$ is a fixed point of the gap operator at that
coupling. Because $I$ is strictly decreasing, $G(\Delta) > G(0) = G_c$,
so every member of this family lies inside the supercritical regime the
conjecture was about, and $G(\Delta)$ increases without bound as
$\Delta$ does. The fixed-point set is a curve in the $(G,\Delta)$ plane
that runs past $\Delta = \Lambda$ and never turns back. Strict
monotonicity of $I$ then converts the conjecture into an exact
threshold: a positive fixed point satisfies $\Delta < \Lambda$ if and
only if
\begin{equation}
  G \;<\; G^{\ast} \;=\; \frac{G_c}{1-\log 2},
  \label{eq:threshold}
\end{equation}
since $\Delta < \Lambda$ is equivalent to $I(\Delta) > I(\Lambda)$,
which is equivalent to $1/(G N_f) > 1/(G^{\ast} N_f)$. The conjecture
was therefore not merely wrong. It was missing exactly one hypothesis,
and the missing hypothesis is identifiable in closed form from
components that were already verified. Equations~(\ref{eq:invertgap})
and~(\ref{eq:threshold}) follow from \texttt{gapIntegral\_pos} and
\texttt{gapIntegral\_strictAnti} together with the closed form; the
library does not currently carry either as a declaration, and they are
presented here as consequences of verified components rather than as
verified statements.

The physical content of Eq.~(\ref{eq:threshold}) is worth stating in
its own right, because it is a result obtained from a failed proof. The
window of couplings in which this model condenses \emph{and} stays
inside its own regime of validity is the bounded interval
$G_c < G < G_c/(1-\log 2)$, a factor of $1/(1-\log 2)$ in coupling, not
the half-line $G > G_c$ that the folklore reading licenses. A
calculation that scans supercritical couplings without that ceiling
will at some point report a condensate the effective description cannot
support, and will do so without any internal signal that it has left
the model's domain.

What the library ships is the surviving statement,
\texttt{gap\_solution\_monotone}
(\texttt{lean/SKEFTHawking/TetradGapEquation.lean}, lines~333--349,
Aristotle run~\texttt{79e07d55}): for couplings $G_1 < G_2$ with
positive fixed points $\Delta_1$ and $\Delta_2$ respectively,
$\Delta_1 < \Delta_2$. Its shape records the refutation as much as its
content does. It quantifies over fixed points supplied as hypotheses,
not over a solution function $\Delta^{\ast}(G)$, so it assumes neither
that a branch has been selected nor that a selected branch stays
anywhere in particular; and it concludes an ordering rather than a
bound, so it says nothing that a large enough $G$ can falsify. The
disproved statement is retained in the module as a commented-out block
at lines~311--325, carrying the counterexample data and the reason the
statement fails, so that a reader of the module encounters the
refutation at the point where the missing theorem would otherwise look
like an oversight.

Downstream, the bound does not silently reappear. Where a later module
needs a condensate that stays below the cutoff, it takes that as an
explicit assumption rather than importing it from the gap equation. The
predicate \texttt{H\_ScalarChannelIsTetradBifurcationOutput}
(\texttt{lean/SKEFTHawking/ScalarRungInterpretation.lean},
lines~170--172) states that a scalar channel arising from the
supercritical branch has condensate VEV bounded by the UV cutoff, and
it is consumed as a hypothesis by
\texttt{mexican\_hat\_vev\_under\_supercritical\_bridge}
(lines~179--184), whose conclusion is exactly as conditional as its
input. The predicate has content rather than being a convenience:
\texttt{bridge\_excludes\_super\_uv\_vev} (lines~190--196) proves the
contrapositive, that a channel whose VEV exceeds the cutoff cannot be
an output of the branch, so a super-cutoff channel refutes the bridge
instead of quietly satisfying it. The assumption is carried in the
program's \texttt{HYPOTHESIS\_REGISTRY} (\texttt{src/core/constants.py})
with an eliminability grade and a named discharge path, namely the
quantitative map from the scalar-channel parameters to the expansion
coefficients of the gap-equation bifurcation, on whose resolution the
predicate becomes a theorem of the gap-equation module. A referee can
see, without reading a proof, exactly which conclusions of the
emergent-gravity sector rest on the sub-cutoff assumption and what
would be required to remove it.

\begin{figure}[h]
\centering
\includegraphics[width=0.85\linewidth]{figures/i1_fig5_gap_counterexample.png}
\caption{Saturation of the gap at the cutoff. Below the critical
coupling the trivial solution $\Delta = 0$ is the only one (dotted
line, \texttt{gap\_trivial\_unique\_subcritical}); above it a positive
fixed point exists (\texttt{gap\_nontrivial\_exists}) and increases
with $G$ (\texttt{gap\_solution\_monotone}). The branch meets the
cutoff at the saturation coupling $G^{\ast} = G_c/(1-\log 2)$ (amber
marker), which is where the counterexample returned in Aristotle
run~\texttt{79e07d55} sits, and exceeds it throughout the amber region,
where the conjectured bound $\Delta < \Lambda$ fails. The onset and the
saturation point are exact; the branch drawn between and beyond them is
an interpolation respecting the square-root onset, not a numerical
solution of Eq.~(\ref{eq:gapeq}).}
\label{fig:gap_counterexample}
\end{figure}

% TODO: the generator fig_i1_gap_counterexample in
% src/core/visualizations.py places the onset at G = 1 while declaring
% the c_4 = Lambda = N_f = 1 convention, in which
% criticalCoupling_formula gives G_c = 2. The rendered figure therefore
% shows a G*/G_c ratio of 6.5 where the verified ratio is
% 1/(1 - log 2) = 3.26. The figure must be regenerated before
% submission; the caption above avoids asserting either numeral.

The general lesson is narrower than ``formal methods find bugs'' and
more useful. A pen-and-paper theorem inherits its hypotheses from the
context in which it is written, and the reader supplies the missing
ones without noticing. A mechanized theorem cannot: the hypothesis list
is finite, explicit, and adversarially searchable, and an automated
prover that cannot close the goal is searching precisely the gap
between the hypotheses that were written and the ones that were meant.
In this case the search terminated in a witness, the witness localized
the omission to a single missing inequality, and the inequality turned
out to have a closed form that bounds the model's regime of validity
from above. None of that required knowing in advance that the
conjecture was false.

\section{Worked Case~3: The Chirality-Wall Axiom Decomposition}
\label{sec:case3}

A proof goal can be out of reach because it is large rather than
because it is deep, and from inside a tactic loop the two failures are
indistinguishable: in both cases nothing closes. The distinction is
operational. A deep goal needs an idea. A large one needs bookkeeping,
and the bookkeeping is mechanical once someone insists on it. The third
worked case is a large goal. The decomposition that made it tractable
also did something no monolithic proof could have done: it forced into
the open a physics assumption that none of the project's assumption
registries recorded.

The theorem is a Standard-Model exclusion, formalized in
\texttt{lean/SKEFTHawking/EWBaryogenesisChiralityWall.lean}. The module
encodes viability of electroweak baryogenesis as a conjunction of two
necessary conditions,
\begin{equation}
\label{eq:ewbgviable}
\texttt{EWBGViable}(a, p, c) \;\equiv\;
\texttt{WallCracked}(a) \;\wedge\;
\mathsf{IsBaryogenesisViable}(p, c),
\end{equation}
where $a \in \mathbb{Z}_{16}$ is the anomaly of the fermion content,
$p$ is a record of finite-temperature effective-potential parameters
(\texttt{EWFiniteTParams} in
\texttt{lean/SKEFTHawking/EWPhaseTransition.lean}), and $c$ is a real
sphaleron-decoupling threshold. \texttt{WallCracked}$(a)$ is $a = 0$:
the anomaly cancels, so the chiral theory admits the gauging step and a
non-perturbative regulator exists. The second conjunct is the
phase-transition condition, $0 < E$ together with
$c\,\lambda\,T_c < E$, where $E$ is the cubic thermal coefficient
\texttt{p.cubic\_coeff} and
$T_c = \sqrt{\mu^2 / c_T}$ is the critical temperature built from the
record's other fields.

The root theorem is
\texttt{sm\_no\_nu\_R\_ewbg\_doubly\_forbidden}. For every parameter
record $p$, every threshold $c$, and under the hypothesis that the
resummed Standard-Model transition is a crossover, it asserts
\emph{both}
$\neg\,\texttt{EWBGViable}(45, p, c)$ and
$\neg\,\texttt{EWBGViable}(48, p, c)$. The two conjuncts are statements
about two different fermion contents and they fail for two different
reasons. Three generations without right-handed neutrinos carry
$3 \times 15 = 45$ Weyl components and $45 \equiv 13 \not\equiv 0
\pmod{16}$, so the wall is intact and the exclusion is independent of
transition order. Adding three right-handed neutrinos gives $48 \equiv
0$, which cracks the wall, and the exclusion then rests entirely on the
crossover. Stating both is not a redundant bundle: neither conjunct
implies the other, and dropping either leaves a fermion content for
which the paper's headline claim is unsupported.

Stated monolithically, that goal is large. Unfolding
Eq.~\eqref{eq:ewbgviable} at both anomaly values expands the negated
conjunction into two decision problems in $\mathbb{Z}_{16}$, two copies
of a three-way real-arithmetic conjunction containing a square root,
and the four bundled positivity fields that \texttt{EWFiniteTParams}
carries with every instance ($0 < \mu^2$, $0 < \lambda$, $0 < c_T$,
$0 \leq E$). De~Morgan over the whole thing needs a case split on
$a = 0$ that the goal does not supply. The pipeline's own guidance
(\S\ref{sec:pipeline}) lists monolithic simplification over goals of
this size as an anti-pattern and sets the target for a submitted
sub-goal at roughly a dozen term-level constraints.

None of the pieces is deep. The $\mathbb{Z}_{16}$ facts are decidable
and close by \texttt{decide}. The crossover exclusion is a two-line
contradiction against a strict inequality. The wall-intact verdict for
the Standard Model is a direct consumption of an upstream arithmetic
theorem. The size came from stating all of them at once, and it is the
only thing that made the goal hard.

% TODO: the monolithic attempt and its failure mode are asserted in the
% project's own working record but have no run identifier: there is no
% ARISTOTLE_THEOREMS entry for EWBaryogenesisChiralityWall and no
% archived run for either the failed monolithic submission or a batch
% closure of the decomposed sub-lemmas. Either a run ID is supplied, or
% this case study stands on the substrate shape alone (as drafted here)
% and no claim about an automated-prover outcome is made.

The module as it ships contains no monolithic statement at all. It
declares sixteen theorems, no \texttt{sorry}, and no project-local
axiom. The longest proof body is eleven tactic lines, in the
biconditional discussed below; no other proof in the module runs past
seven. The root theorem's own proof is a two-element anonymous
constructor. Reading upward from it, the chain that discharges the
exclusion is:

\begin{itemize}
\item \texttt{sm\_no\_nu\_R\_wall\_intact}: $45 \neq 0$ in
$\mathbb{Z}_{16}$, consuming
\texttt{Z16AnomalyComputation.three\_gen\_anomalous}.
\item \texttt{ewbg\_forbidden\_if\_wall\_intact}: an intact wall
refutes the first conjunct of Eq.~\eqref{eq:ewbgviable} for any $p$ and
any $c$.
\item \texttt{sm\_no\_nu\_R\_ewbg\_blocked}: the two above, composed at
$a = 45$.
\item \texttt{ewbg\_forbidden\_if\_transition\_crossover}: a crossover
refutes the second conjunct for any anomaly, invoking
\texttt{EWPhaseTransition.crossover\_excludes\_baryogenesis} in its
proof body.
\item \texttt{sm\_with\_3nu\_R\_ewbg\_forbidden\_under\_klrs}: the
previous item at $a = 48$, under the tracked crossover hypothesis.
\end{itemize}

Around that spine sit the structural lemmas that make the
decomposition auditable rather than merely shorter:
\texttt{wall\_intact\_xor\_cracked} and
\texttt{wall\_intact\_or\_cracked} (disjointness and exhaustivity of the
wall predicate over $\mathbb{Z}_{16}$);
\texttt{sm\_no\_nu\_R\_anomaly\_eq\_neg\_three}, which identifies the
$45$-component content with the canonical $-3$ representative used in
the anomaly classification;
\texttt{sm\_single\_gen\_with\_nu\_R\_wall\_cracks} and
\texttt{sm\_with\_3nu\_R\_wall\_cracks} at one and three generations;
\texttt{sm\_with\_3nu\_R\_ewbg\_iff\_first\_order\_strong}, which
records that with the wall cracked the compound predicate collapses to
the phase-transition condition alone; and the quantitative anchor
\texttt{sm\_klrs\_overshoot\_ratio\_gt\_threshold}.

% TODO: figures/i1_fig6_chirality_wall_tree.png must be regenerated from
% fig_i1_chirality_wall_tree in src/core/visualizations.py before this
% caption is accurate. The rendered leaf labels are paraphrases
% ("Sublemma A.1") rather than declaration names, and the rendered
% subtitle asserts a single-batch Aristotle closure for which no run
% record exists. The tree below is the one the module actually has.
\begin{figure}[h]
\centering
\includegraphics[width=0.95\linewidth]{figures/i1_fig6_chirality_wall_tree.png}
\caption{Decomposition of the chirality-wall exclusion. The root
\texttt{sm\_no\_nu\_R\_ewbg\_doubly\_forbidden} (carmine) is a
conjunction over two fermion contents whose unfolded form spans two
$\mathbb{Z}_{16}$ decision problems, two copies of the
phase-transition condition, and the four bundled positivity fields of
\texttt{EWFiniteTParams}. It splits into two independent obstructions
(amber): an intact chirality wall at anomaly $45$, which excludes
baryogenesis for any transition order, and a crossover transition at
anomaly $48$, which excludes it for any wall state. Each branch
resolves into named module lemmas (steel blue), none of whose proof
bodies exceeds eleven tactic lines. The biconditional
\texttt{ewbg\_forbidden\_iff\_wall\_intact\_or\_not\_viable} is the
leaf that made the sphaleron-rate dependency visible.}
\label{fig:chirality_wall_tree}
\end{figure}

The payload of this case is not that the decomposition worked. It is
what the decomposition exposed. One of the sub-lemmas,
\texttt{ewbg\_forbidden\_iff\_wall\_intact\_or\_not\_viable}, is a
biconditional rather than an implication: it asserts that every failure
of \texttt{EWBGViable} is attributable to the wall, to the transition,
or to both, and that there is no third failure mode. Writing that
statement out is what raised the question the monolith never had to
answer. Exhaustivity holds only relative to what the compound predicate
records, and the sphaleron rate, which is the physical quantity the
whole decoupling criterion is about, appears in
Eq.~\eqref{eq:ewbgviable} nowhere except as the free threshold
parameter $c$. Nothing in the predicate constrains the rate itself. The
biconditional is therefore true as stated only under a further
assumption about the sphaleron rate, and that assumption had no entry
in any of the project's assumption registries.

The audit closed it by proving it rather than by declaring it. The
module defines the dimensionless suppression factor
$\Sigma(v,T) = e^{-v/T}$ explicitly, as a structural stand-in whose
limits are the physically meaningful content ($\Sigma \to 0$ deep in
the broken phase, $\Sigma \to 1$ in the symmetric phase), and then
bounds it from both sides:
\texttt{sphaleronSuppression\_pos} gives $0 < \Sigma(v,T)$ at any
finite $(v,T)$, and
\texttt{sphaleronSuppression\_le\_one} gives $\Sigma(v,T) \leq 1$
whenever $v/T \geq 0$. A factor that is strictly positive and never
exceeds unity cannot open a failure channel of its own, so the rate
enters the exclusion only through the threshold, which is exactly what
the biconditional's exhaustivity requires. The threshold itself stays a
free real parameter in every statement in the module, including the
root theorem; its numerical value lives in the project's
experimental-parameter table and never inside a theorem. Consistent
with that resolution, no sphaleron-rate hypothesis appears today in
\texttt{AXIOM\_METADATA}, in \texttt{PLACEHOLDER\_THEOREMS}, in
\texttt{HYPOTHESIS\_REGISTRY} or in \texttt{BORDISM\_HYPOTHESES}: there
is nothing left to record, because the assumption became two theorems.
The module ships zero new axioms.

The general point is that an implicit dependency is invisible in
proportion to the size of the statement that hides it. A monolithic
proof of the root theorem would have been correct, would have passed
every mechanical gate the pipeline applies, and would have carried the
same unexamined assumption indefinitely, because no gate in the pipeline
can ask a question about an assumption nobody has named. Decomposition
is what turns an implicit assumption into a statement, and a statement
is something a reviewer can refuse.

What the root theorem does still assume, it assumes visibly. Exactly
one tracked hypothesis appears in its type: \texttt{H\_KLRS\_SM\_Crossover}~$p$,
the proposition that the fully resummed Standard-Model transition is a
crossover at the physical Higgs mass. This is an external lattice
result and is not derived here. It is recorded in
\texttt{HYPOTHESIS\_REGISTRY} (\texttt{src/core/constants.py}), graded as
an external-boundary input of hard eliminability and naming this
theorem among its dependent theorems, and its evidence is quantitative:
the four-dimensional lattice analysis of the SU(2)-Higgs model places
the endpoint of the first-order line at $m_H = 66.5 \pm 1.4$~GeV and
transforms it perturbatively to $m_H = 72.4 \pm 1.7$~GeV for the full
Standard Model, above which only a rapid crossover is
seen~\cite{CsikorFodorHeitger1999}; the earlier three-dimensional
analysis had bracketed the endpoint between $70$ and $95$~GeV, most
likely near $80$~GeV~\cite{KLRS1996}. Against the observed $m_H =
125.20$~GeV the margin is wide, and the module states it as a
falsifiable arithmetic claim rather than a remark:
\texttt{sm\_klrs\_overshoot\_ratio\_gt\_threshold} proves
$1.5 < 125.20 / 72.4$ by \texttt{norm\_num}.

Three registries carry the project's assumption inventory, and which
one an assumption lands in is itself a claim that has to be checked.
\texttt{AXIOM\_METADATA} holds project-local axioms and their retirement
record, graded by eliminability; \texttt{HYPOTHESIS\_REGISTRY} holds
tracked propositions that consumers take as explicit hypothesis
parameters; \texttt{PLACEHOLDER\_THEOREMS} holds declarations proved
\texttt{True := trivial}, which Invariant~9 forbids any proof, formula
or paper claim from leaning on; \texttt{BORDISM\_HYPOTHESES} holds
bordism-theoretic inputs deferred to future Mathlib infrastructure. The
grades are not static. The entry \texttt{gapped\_interface\_axiom} was a
live project-local axiom and is now graded \texttt{closed}, because the
axiom was replaced by a tracked proposition (\texttt{TPFConjecture} in
\texttt{lean/SKEFTHawking/SPTClassification.lean}) that every consumer
takes as an explicit hypothesis parameter. The substantive conjecture is
unchanged; what changed is that it is now visible in four type
signatures instead of silently present in every closure below the
axiom. That is the same move as the sphaleron case one step further
along: name the assumption, then discharge it if you can.

\section{The Fourteen-Stage Wave Execution Pipeline}
\label{sec:pipeline}

The pipeline is organized as fourteen sequential stages, each with
a gate that must pass before the next stage begins. Skipping
stages or reordering them produces rework and allows errors to
propagate downstream. Wave-level work (adding new physics,
extending formalization, or producing experimental predictions)
proceeds through all fourteen.

\begin{figure}[h]
\centering
\includegraphics[width=0.95\linewidth]{figures/i1_fig2_pipeline_14_stages.png}
\caption{Fourteen-stage wave-execution pipeline. Each cell shows
the stage number, short name, and gate condition. New stages
(3a Lean MCP loop, 13 Adversarial review, 14 Meta-process QI,
amber) supplement the original eleven (steel blue). Stage 3a
replaces the prior write-tactic / \texttt{lake build} / parse /
iterate loop ($\sim$15 s per cycle) with millisecond-scale
LSP-state inspection and \texttt{lean\_multi\_attempt} candidate
testing. Stage 13 is the fresh-context adversarial-reviewer
pattern. Stage 14 captures meta-process improvements as
project-level memory.}
\label{fig:pipeline}
\end{figure}

\begin{table}[h]
\centering
\small
% Hand-authored for paper I1; tracks docs/WAVE_EXECUTION_PIPELINE.md.
% This is the 14-stage pipeline that supersedes paper15's 12-stage table.
% If WAVE_EXECUTION_PIPELINE.md is updated, sync this table.
\begin{tabular}{cll}
\hline
\textbf{Stage} & \textbf{Name} & \textbf{Gate} \\
\hline
1   & Constants \& parameters       & imports succeed; provenance verified \\
2   & Formulas                       & every function has Lean ref \\
3a  & Lean --- interactive MCP loop  & proof closes (or decomposed to sub-lemma stubs) \\
3b  & Lean --- sorry registration    & lake build compiles, SORRY\_GAPS accurate \\
4   & Aristotle (fallback)           & residual sorrys filled, lake build clean \\
5   & Lean build verification        & zero sorry, counts match \\
6   & Python tests                   & all tests pass \\
7   & Cross-layer validation         & validate.py all checks pass \\
8   & Visualizations                 & all PNGs generated \\
9   & Figure review                  & LLM review all PASS \\
10  & Paper draft                    & paper\_provenance check passes \\
11  & Notebooks                      & notebook\_exec + viz\_consistency pass \\
12  & Document sync                  & full validate.py passes, counts consistent \\
13  & Adversarial review             & fresh-context Opus sweep shows zero BLOCKERs \\
14  & Meta-process QI                & systemic findings logged to QI register (advisory) \\
\hline
\end{tabular}

\caption{Pipeline stages with their gates, autogenerated from
\texttt{WAVE\_EXECUTION\_PIPELINE.md}.}
\label{tab:stages}
\end{table}

The current pipeline extends the original twelve-stage version
(paper15 draft, April 2026) with three additional stages.
Stage~3a, interactive Lean development via the language-server
protocol, replaces the previous ``write tactic, run \texttt{lake
build}, parse error, iterate'' development loop (approximately
fifteen seconds per cycle) with a millisecond-latency interactive
loop using the \texttt{lean-lsp-mcp} tool integration. The
interactive loop reads the proof state at any line and column,
tests candidate tactic sequences without editing the file, and
returns ``no goals'' when the proof is complete. Per tactic-test cycle, the interactive
loop replaces a $\sim\!15$\,s lake-build round-trip with
millisecond-scale LSP state inspection. End-to-end
proof-construction time on a given theorem improves by a factor
determined by the number of tactic-test cycles the proof requires;
we have not measured this factor on a controlled benchmark set, and
quantitative speedup claims would require such a benchmark to be
load-bearing.

Stage~13, fresh-context adversarial review, runs a Lean-aware
adversarial reviewer agent in a fresh context window against each
paper or bundle near submission readiness. The agent works through
eight finding classes (citation integrity, parameter provenance,
formal-verification-claim substantiation, cross-paper
consistency, narrative grounding, assumption disclosure, count
freshness, and production-run health, the same eight classes
enumerated in detail in Section~\ref{sec:adversarial}) and emits
a structured markdown report.
Findings are auto-picked up as graph nodes; ``BLOCKER'' findings
flip the affected readiness gate to blocked. Stage~13 is meaningful
only on a codebase that already passes its own internal validation
(Stages~1--12); it backstops the failure classes that internal
validation cannot detect by construction.

Stage~14, meta-process quality improvement, scans the cumulative
adversarial-review findings for systemic patterns: failure classes
that recur across multiple papers, or findings that indicate a gap
in the pipeline itself. Such patterns are recorded in a quality-
improvement register with owners and target dates. Stage~14 is
advisory; it never blocks publication. Its purpose is feeding
pipeline improvements back into subsequent waves. The April~2026
external-adversarial-review round, which surfaced a thirteen-
dimensional problem space across the project, was the seed event
for the register.

Stages~1--12 are the substantive computational and formalization
work. Stage~1 establishes the canonical sources of truth: physical
constants in \texttt{constants.py}, formulas in
\texttt{formulas.py}, figure functions in \texttt{visualizations.py}.
Stage~2 implements formulas with provenance metadata. Stages~3--4
formalize and prove (interactive plus automated). Stage~5 verifies
the formal layer compiles cleanly. Stage~6 runs Python tests.
Stage~7 runs the automated cross-layer validation suite
(enumerated live by \texttt{validate.py --list}; the count is
deliberately not quoted here, because a census figure written into
prose is stale the moment a check is added: this sentence
read ``thirty-three'' against a live seventy-nine)
spanning formula--theorem correspondence, parameter provenance,
count freshness, citation integrity, and bundle consistency
(Sections~\ref{sec:provenance} and~\ref{sec:adversarial}).
Stage~8 generates publication-
quality figures. Stage~9 runs a figure-reviewer agent. Stage~10
drafts the paper with full provenance for every claim. Stage~11
produces reproducible computational notebooks. Stage~12 syncs
documentation and runs counts regeneration. Each stage's gate must
pass before proceeding.

The pipeline's invariants must hold at all times, not just at stage
completion:

\begin{enumerate}
\item \textbf{Canonical formulas.} \texttt{formulas.py} is the only
place physics formulas live.

\item \textbf{Canonical constants.} \texttt{constants.py} is the only
place experimental parameters and the theorem registry live.

\item \textbf{Canonical visualizations.} \texttt{visualizations.py}
is the only place figure functions live.

\item \textbf{Every formula has a Lean theorem.} Every Lean theorem
has a proof (zero sorry target).

\item \textbf{Every computed quantity has bounds.} Physical bounds
are tested and enforced by validation.

\item \textbf{Every paper claim traces to computation.} Numerical
claims match formula output within 0.5\%.

\item \textbf{Narrative derives from data.} Feasibility claims need
computed support.

\item \textbf{Every parameter has verified provenance.} Each value
traces to a published source via the
\texttt{PARAMETER\_PROVENANCE} table.

\item \textbf{Placeholder theorems are non-load-bearing.} Theorems
proved as \texttt{True := trivial} are documentation markers
only. Tracked in a \texttt{PLACEHOLDER\_THEOREMS} registry; not
counted as substantive.

\item \textbf{No heartbeat overrides in proof bodies.} A proof body
that hits the heartbeat limit is evidence that the proof
architecture is wrong; the fix is sub-lemma decomposition, not
raising the budget.

\item \textbf{Every external bibitem has a primary-source cache
file.} Each non-in-prep entry in the citation registry carries a
resolvable path to a cached copy of the cited primary source,
held in a per-phase literature cache maintained alongside the
repository, structurally preventing the hallucinated-citation
failure mode.

\item \textbf{Provenance DOIs resolve to the registry.} Every DOI
mentioned in any parameter-provenance source must resolve to a
registry bibkey; every cited bibkey must exist in the registry.

\item \textbf{Stage~14 quality-improvement register is auto-
regenerated and manually curated, never wiped.} Closed items
are preserved verbatim across regenerations.

\item \textbf{Every paper-shaped output lifts into a publication
bundle.} New drafts, section additions, and notebook companions
identify their target bundle at Stage~1 and record it in the
draft-mapping document at Stage~12. No new stand-alone paper
drafts are created without explicit user authorization to add a
new bundle target beyond the current roster.

\item \textbf{Every new project-local axiom requires explicit
user sign-off.} Locked in May~2026. No axiom ships on
recommendation-only authority; every new axiom carries either a
discharge plan (a scheduled wave producing a substantive
constructive proof) or a documented argument that no constructive
proof is feasible against the current upstream mathematics
library. The posture is that axioms are temporary scaffolding,
not permanent commitments: the project's last pre-existing
project-local axiom was retired into a tracked-hypothesis
\texttt{Prop} on 2026-05-19 (Section~\ref{sec:case3}), taking the
live project-local axiom count to zero.

\end{enumerate}

Invariant~9 addresses a subtle inflation problem:
\placeholdertheorems{} of our \totaltheorems{} theorems are
\texttt{True := trivial} placeholders that serve as documentation
markers. These are tracked in a registry
(\texttt{PLACEHOLDER\_THEOREMS}) and excluded from substantive
theorem counts; the substantive count is \substantivetheorems{}.

Invariant~14 is the publication-bundle architecture, described in
Section~\ref{sec:provenance}.

%% ─────────────────────────────────────────────────────────────────
%% §7 Preemptive-strengthening discipline
%% ─────────────────────────────────────────────────────────────────

\section{The Preemptive-Strengthening Discipline}
\label{sec:strengthening}

The strengthening pass is an audit that runs at the end of each
wave, examining each new theorem statement for six recurring
anti-patterns. Originally retroactive (the audit ran after
theorems were stated and proved, with retroactive cuts applied
where the audit identified a vacuous statement), the pass has
been moved to the prospective phase: each theorem statement is
checked against the same six anti-patterns \emph{before} the
statement is written. The reduction in audit cost is
substantial: the wave that introduced the discipline (Phase~6c.3)
required twelve retroactive cuts; the next wave (Phase~6b.1)
required five; the wave after that (Phase~6m R6) required zero.

The six anti-patterns are:

\begin{itemize}

\item \textbf{P2: Bundle redundancy.} A theorem stating a
conjunction is redundant if dropping any conjunct produces a
statement of equal information content. The pattern occurs when
multiple conjuncts are algebraically equivalent: for instance,
$\mathrm{satisfiesAt}\ m\ \mathrm{WEP}\ \wedge\ \mathrm{satisfiesAt}\
m\ \mathrm{EEP}\ \wedge\ \mathrm{satisfiesAt}\ m\ \mathrm{SEP}$ when
$\mathrm{violationLevel}\ m = \mathrm{none}$ is one substantive
theorem dressed as three. The fix is to write the substantive form
($\mathrm{violationLevel}\ m = \mathrm{none}$) plus a separate
extraction lemma if downstream callers need the bundle shape.

\item \textbf{P3: Multiplication-only triviality.} A theorem
whose substantive content is exhausted by an obvious arithmetic
identity ($0.34 = 2 \times 0.17$, for instance) is structural, not
physical. The fix is to replace with a load-bearing comparison
that separates a conformant case from a violator
($0.34 < 1.0$).

\item \textbf{P4: Self-discharging axiom.} An axiom whose
statement reduces to truth via algebra alone is vacuous. The fix
is to either remove the axiom or write the substantive content
elsewhere.

\item \textbf{P5: Falsifier as restated hypothesis.} A
``falsifier'' theorem whose body restates its hypotheses without
exercising any computation is structural-tautology repackaging.
The fix is to ensure the falsifier exercises the hypothesis.

\item \textbf{P6: Cross-module bridge integrity.} If a
theorem's docstring references another module's theorem, the body
must \texttt{import} that module and \emph{call} the theorem,
not merely list it as a documentation pointer. The audit pattern
is: every cross-reference in a docstring must be backed by a Lean
call in the body. A bridge whose body never invokes the named
target is documentation-only and rots silently as the codebase
evolves.

\item \textbf{P7: Defining the conclusion.} If a function is
defined as \texttt{:= <obvious target>}, a theorem that asserts
$0 \leq C/A$ for that function is trivially true. The fix is to
either push the substantive content elsewhere (a tracked-hypothesis
proposition, or a project-local infrastructure module) or commit to
a non-trivial concrete model.

\end{itemize}

The discipline is applied as a five-question prospective checklist
on each theorem statement before it is written. The reduction in
strengthening-pass cost is the primary metric, but a ruthless
post-wave review remains mandatory: empirically, the
prospective discipline catches the more obvious failure modes
($\exists$-absorption, biconditional-tautology in conditional
violators) but misses subtler patterns: identity-function wrappers
($\theta\ h_a = h_b$), within-own-$\pm 2\sigma$-band tautologies,
pairwise-distinctness on inductive constructors, and definitional-
unfolding-as-physics. Both passes are mandatory.

The discipline is not aesthetic. Each anti-pattern corresponds to
a specific failure of the theorem-statement-as-falsifiable-claim
property: a theorem that can be discharged by \texttt{rfl} or
\texttt{decide} or \texttt{not\_lt.mpr h\_disagree} or
$0 \leq C$ for any positive $C$ does not constitute a falsifiable
physical claim. The strengthening discipline ensures that what we
prove is what we claim.

%% ─────────────────────────────────────────────────────────────────
%% §8 Sentence-level provenance + cross-paper consistency clusters
%% ─────────────────────────────────────────────────────────────────

\section{Sentence-Level Provenance and Cross-Paper Consistency Clusters}
\label{sec:provenance}

\begin{figure}[h]
\centering
\includegraphics[width=0.95\linewidth]{figures/i1_fig3_sentence_state_clusters.png}
\caption{Sentence-state claim clusters across paper bundles
(schematic). Three papers (rectangles) each contain six sentences
(circles). Curved lines connect sentences sharing a
\texttt{claim\_cluster}. The amber cluster~C2 has
\texttt{cross\_bundle: true} and links sentences across all three
papers; clusters C1 (steel blue) and C3 (sage) are
single-bundle. Cross-bundle clusters survive bundle
consolidation, providing the per-bundle Stage-13
adversarial-reviewer with a stable provenance index even when
sentences are lifted into different paper bundles. Walked by
\texttt{validate.py --check bundle\_consistency}.}
\label{fig:sentence_clusters}
\end{figure}

% TODO: substantive draft. Substrate to draw from:
%  - prose_state.json schema in scripts/sentence_state.py
%  - cross-paper consistency cluster mechanism in
%    scripts/cluster_detect.py + scripts/bundle_clusters.py
%  - claim_clusters.json + cluster_bundle_index.json semantics
%  - validate.py CHECK 21 (bundle_consistency) walks the registry
%    and flags cross-bundle drift
%  - The publication-bundle architecture per docs/PAPER_STRATEGY.md
%  - The bundle_metadata.json schema (Phase 7a sub-wave 7a.1.1) and
%    the BUNDLE_LIFT_PROCEDURE.md canonical lift workflow
%
% Section thesis: per-paper provenance tracks WHERE each numerical
% claim comes from; sentence-level provenance tracks WHAT each
% prose sentence asserts and which sentences across the program
% make the same assertion. Cross-paper consistency clusters are
% the mechanism that prevents the same fact from being stated
% differently in two papers --- a publication-readiness invariant
% that emerges when the program ships >10 papers.

Sentence-level provenance is the mechanism that scales the per-claim
audit discipline from a single paper to a multi-paper publication
program. Each prose sentence in every paper draft acquires a record
in a structured registry (\texttt{papers/<paper>/prose\_state.json},
schema-defined in \texttt{scripts/sentence\_state.py}) that captures
the claim type (numeric, theorem-reference, citation, parameter,
formal-claim, qualitative, transition, methodology, or metaclaim), a
chain of backing links (formula in \texttt{formulas.py}, theorem in
the Lean source, axiom in \texttt{AXIOM\_METADATA}, parameter in
\texttt{PARAMETER\_PROVENANCE}, citation in
\texttt{CITATION\_REGISTRY}, hypothesis in
\texttt{HYPOTHESIS\_REGISTRY}), an agent verdict (PASS, FAIL, WARN,
INFO, UNGROUNDED, TRANSITION) from the most recent claims-reviewer
run, and a content-hash sentence identifier
(\texttt{sentence:<paper>:<section\_slug>:<sha8>}) that survives
prose edits within a paper but flips when the substantive claim
changes.

The sentence ID is the stable handle. Adversarial-review findings
are pinned to sentence IDs rather than line numbers, so an author's
prose edit in one paragraph does not invalidate findings on
neighboring paragraphs. Re-review is incremental: only sentences
whose hash has changed since the last claims-reviewer run need fresh
chain resolution; the rest are loaded from the prior
\texttt{claims\_review.json}.

Cross-paper consistency clusters are the next layer. A claim made in
two or more papers (a numerical bound, a theorem statement, a
parameter value) groups its sentence IDs into a cluster registered
in \texttt{papers/claim\_clusters.json}, with a
\texttt{cross\_bundle: true} flag for clusters spanning two or more
publication bundles. The validation check
\texttt{validate.py --check bundle\_consistency} walks
the registry and flags any cluster whose member sentences disagree
on the cluster's load-bearing fact. The load-bearing example is
GW170817 vs vestigial-graviton: the L1 PRL splash and the D3 deep
paper both report the LIGO $\Delta c/c$ bound and the
$7\times 10^{14}$ falsification factor; the
\texttt{bundle\_consistency} check enforces that the
two reports agree character-for-character on the numerical content.

The publication-bundle architecture (codified in the process
document \texttt{docs/PAPER\_STRATEGY.md} in the public repository)
consolidates the project's seventy-plus mapped sources (per-wave
drafts plus Lean-only synthesis handles) into the
publication targets that frame enumerates organized into five tiers: a flagship review
(\texttt{F}), eight themed deep papers
(\texttt{D1}--\texttt{D8}), three PRL-style headline letters
(\texttt{L1}--\texttt{L3}), three infrastructure papers (\texttt{I1},
\texttt{I2}, \texttt{I3}), and two experimental letters (\texttt{E1},
\texttt{E2}). Each per-wave draft maps to one bundle target via an
append-only assignment table
(repository document \texttt{docs/PAPER\_DRAFT\_MAPPING.md}); each
bundle accumulates
contributions from its sources via a documented lift workflow
(repository document \texttt{docs/BUNDLE\_LIFT\_PROCEDURE.md},
fourteen canonical steps);
late-arriving Phase~6X waves absorb into already-drafted bundles via
a hot-swap protocol
(repository document \texttt{docs/LATE\_PHASE6\_ABSORPTION\_PROTOCOL.md},
stages~A--G with explicit branches D.1 / D.2 / D.3 / D.4 by bundle
state). Per-bundle
state is captured in \texttt{papers/<bundle>/bundle\_metadata.json},
whose fields include the three review-stage
statuses (Stage~9 / 10 / 13), a freshness flag, the source-manifest
last-regeneration timestamp, the stage-13 review document path, and
the recorded-review timestamps consumed by the readiness semantics
of Section~\ref{sec:adversarial}. Pipeline
Invariant~14 codifies the contract: every paper-shaped output lifts
into a bundle target identified at Stage~1 and recorded at
Stage~12. The architecture is what makes many-source-to-few-publication
consolidation tractable without sacrificing
per-claim provenance: the sentence-level audit trail moves with the
prose; the cluster registry preserves cross-bundle consistency; the
bundle metadata captures every reviewer agent's verdict.

Two deposit-time process standards extend the bundle architecture
to publication mechanics. First, the manuscript-level disclosure
of machine assistance is \emph{register-derived rather than
prose-discretionary}: a single canonical disclosure text
(repository document \texttt{docs/DISCLOSURE\_TEXT.md}) exists in
exactly two variants, differing only in one conditional sentence
crediting the Aristotle automated theorem prover, and the variant
each bundle carries is computed deterministically
(\texttt{scripts/aristotle\_usage\_by\_bundle.py}) from the
canonical registers (the Aristotle run registry, the per-paper
declared Lean modules, the draft-to-bundle mapping, and the Lean
dependency export) rather than chosen by the author. Second,
deposit sequencing (repository document
\texttt{docs/ARXIV\_DEPOSIT\_PLAN.md}) is gated on completion of
the content-level attribution sweep described in
Section~\ref{sec:adversarial}, and the deposit-time conversion of
in-preparation self-citations into deposited identifiers is a
single documented scripted pass
(\texttt{scripts/convert\_inprep\_citations.py}) that updates the
citation registry, rewrites the affected bibliography items, and
writes the cache stubs the primary-source invariant requires,
keeping Invariant~11 unbroken across the deposit transition.

%% ─────────────────────────────────────────────────────────────────
%% §9 Three-layer architecture extended (Aristotle integration)
%% ─────────────────────────────────────────────────────────────────

\section{Aristotle Automated Theorem Prover Integration}
\label{sec:three_layer_extended}

% TODO: substantive draft. Substrate to draw from:
%  - paper15 §5 (Aristotle integration)
%  - The priority batching protocol
%  - Run 78dcc5f4 (25 theorems in a single session)
%  - Run 79e07d55 (G_c critical coupling)
%  - The interactive MCP loop as Stage 3a (replaces Stage 4 fallback
%    workflow described in paper15)
%
% Section thesis: Aristotle is now Stage 4 --- a fallback for
% sorries that survive the interactive Stage 3a. The interactive
% loop is roughly 1000x faster for proof construction; Aristotle
% is reserved for sorries that genuinely resist interactive proof
% with sub-lemma decomposition.

The role of automated theorem proving in the pipeline has shifted
substantially over the year of operation. The original twelve-stage version of the pipeline (an earlier
internal draft, April 2026) treated Aristotle as the primary
proof-construction tool: a Lean
theorem statement was registered with a strategy-hint comment in
its docstring, the entire project was submitted in a priority
batch, and the returned proofs were integrated after a manual
diff review. The cycle time was approximately one batch per day,
with batch sizes of ten to forty theorems.

The introduction of Stage~3a (interactive Lean development via
the \texttt{lean-lsp-mcp} integration) changed the balance.
Interactive iteration on a single theorem now takes seconds, not
hours; the tool reads the goal state at the cursor, tests four to
six candidate tactic sequences without editing the file, and
returns the resulting goal state for each. A typical proof that
would have been submitted to Aristotle in the original workflow
now closes interactively in three to five tactic-test cycles
(see Section~\ref{sec:pipeline} for the per-cycle latency
discussion). Aristotle has correspondingly retreated from
default proof-construction tool to fallback-for-residual-sorries.

The residual role is genuine, not vestigial. Three classes of
proof retain a strong indication for Aristotle. The first is
proofs whose tactic search space is too large for interactive
candidate-testing: lengthy algebraic manipulations on
non-commutative ring expressions, multi-step bidegree
decompositions with case splits at every step, and Sobolev-type
inequalities with auxiliary functions whose definitions are not
suggested by the local goal state. The second is proofs that
submit cleanly only after a structural decomposition into
sub-lemmas (the chirality-wall case of Section~\ref{sec:case3}),
where the sub-lemmas are individually small but numerous enough
that interactive proof of all of them is more expensive than a
single batch submission. The third is the diagnostic role
illustrated in Section~\ref{sec:case1} and
Section~\ref{sec:case2}: an Aristotle failure on a mathematically
plausible statement is informative when the failure is reported
as an explicit counterexample, surfacing a wrong axiom or a
folklore-but-false claim that interactive iteration would not
have surfaced as efficiently.

The project's Aristotle activity to date is summarized by
\aristotleproved{} cumulative theorems closed across
\aristotleruns{} priority-batch runs, registered in the
\texttt{ARISTOTLE\_THEOREMS} table in
\texttt{src/core/constants.py}. Each entry records the theorem's
fully qualified name, the run UUID, the priority tier (1 highest,
3 lowest), and the integration date. The single-batch record is
twenty-five theorems closed in run \texttt{78dcc5f4}, a
Phase~5c Wave~1--7 multi-domain batch spanning S-matrix
unitarity and Verlinde formulas at $k = 1, 2$ in
\texttt{SU2kSMatrix}; the modular-tensor-category fingerprint
in \texttt{RibbonCategory}; the Rokhlin invariant in
\texttt{SpinBordism}; the $E_8$-lattice positivity proof in
\texttt{E8Lattice}; restricted-quantum-group corepresentations
in \texttt{RestrictedUq}; the integrated-autocorrelation-time
estimator bounds in \texttt{VerifiedJackknife}; and seven public
Hopf-algebra axioms in \texttt{Uqsl2Hopf}. The operationally
most significant batch is the Phase~5d gap-equation batch
\texttt{79e07d55}, which both produced the
gap-solution-bounded counterexample of Section~\ref{sec:case2}
and closed nine companion gap-equation theorems in the same
submission.

The Aristotle service runs on a fixed Lean toolchain
(\texttt{leanprover/lean4:v4.28.0}) with a corresponding pinned
Mathlib commit, while the project's local toolchain is
currently \texttt{leanprover/lean4:v\leantoolchain} with a separately
pinned Mathlib commit (see
\texttt{lean/lean-toolchain} and \texttt{lean/lakefile.toml}).
Aristotle-returned proofs that exercise a 4.29.0-only feature
in a tactic body fail at integration time on the project's
toolchain; we run a \texttt{lake build} verification on every
integrated batch (Stage~5 of the pipeline) before
\texttt{ARISTOTLE\_THEOREMS} entries are added. The toolchain
skew is the operational reason the project uses Aristotle
\emph{sparingly}: every batch incurs a non-zero
toolchain-compatibility audit cost, and when the cost dominates
the proof-construction cost, the interactive Stage~3a loop is
the better tool. All worked-case theorems referenced in
Sections~\ref{sec:case1}--\ref{sec:case3} have been verified to
compile cleanly on the project's local \leantoolchain{} toolchain
post-integration. The integration protocol is unchanged across the
shift: never \texttt{--integrate} on submission, always retrieve
first, review the diff archived in the repository under
\texttt{docs/aristotle\_results/run\_<timestamp>/diff.patch}, and
integrate only after axiom-hygiene verification (Invariant~10:
no heartbeat overrides introduced; new axioms registered in
\texttt{AXIOM\_METADATA}; \texttt{lake build} clean post-merge).

%% ─────────────────────────────────────────────────────────────────
%% §10 Stage 13 fresh-context adversarial reviewer
%% ─────────────────────────────────────────────────────────────────

\section{The Fresh-Context Adversarial-Reviewer Pattern}
\label{sec:adversarial}

% TODO: substantive draft. Substrate to draw from:
%  - .claude/plugins/physics-qa/agents/adversarial-reviewer.md
%  - WAVE_EXECUTION_PIPELINE.md Stage 13 section
%  - The April 2026 external-adversarial-review round that
%    surfaced the 13-dimension problem space (seed event)
%  - The 8 finding classes (citation integrity, parameter
%    provenance, narrative grounding, cross-paper consistency,
%    count freshness, formal-verification-claim substantiation,
%    assumption disclosure, review-state hygiene)
%  - The supersession ledger semantics (docs/review_finding_supersessions.json)
%
% Section thesis: the adversarial reviewer runs in a fresh context
% (no memory of the author's prior justifications), works through
% 8 finding classes systematically, and produces a structured
% markdown report that flips readiness gates to blocked when
% BLOCKER findings are present. The fresh-context property is
% load-bearing --- the reviewer must not be persuadable by the
% author's prior context.

The fresh-context adversarial reviewer is the pipeline's
load-bearing safeguard against the failure classes that internal
validation cannot detect by construction. The April~2026
external-adversarial-review round on the project's per-paper drafts
exposed a thirteen-dimensional problem space (thirteen distinct
failure-mode categories that none of the project's internal
\texttt{validate.py} checks were designed to surface), including
wrong-target arXiv citations (one paper cited a graph-neural-network
arXiv ID where the intended citation was a torsion-cosmology
paper), parameter values that had drifted from their primary
sources after late-stage rounding, narrative overclaims that were
factually wrong (a ``first machine-checked'' claim that proved to
have an unrecognized prior art in Coq), and ``formally verified''
prose claims whose Lean theorems were placeholder
\texttt{True := trivial} markers. The seed event motivated
Stage~13 as a permanent pipeline component.

The agent's design rests on the \emph{fresh-context} property: the
reviewer runs in a context window that has not seen the author's
prior justifications, internal communications, or in-development
state. By construction, a reviewer in the same context as the author
has internalized the author's framing; the reviewer's findings are
limited to what the author's framing permits. A fresh-context
reviewer reads the paper as a reader would, with no privileged
backchannel into the author's intent, and emits findings against
the artifact rather than against the author. The architectural
choice is that the agent runs in a separate Claude Code session
(separate context window, separate skill loading); the prompt is
the \texttt{physics-qa} plugin's adversarial-reviewer agent
definition plus the paper or bundle target.

The agent works through eight finding classes in strict order, one
class at a time, against eleven readiness gates:
(1)~CitationIntegrity: every \texttt{\textbackslash bibitem} is
fetched against arXiv or DOI and compared character-for-character;
(2)~ParameterProvenance: every numerical parameter is checked
against \texttt{PARAMETER\_PROVENANCE}; (3)~LeanProofSubstance:
every cited Lean theorem is read at its declaration site and its
proof body inspected for placeholder patterns; (4)~CrossPaperConsistency:
the paper is diff'd against every other paper for shared
constructs; (5)~NarrativeGrounding: first-claims, unification
claims, attribution claims, feasibility claims, and
simulation-evidence claims are scanned against the artifact;
(6)~AssumptionDisclosure: every cited theorem's load-bearing
hypotheses must appear in prose; (7)~CountFreshness: count
literals must match the canonical \texttt{counts.tex} macros;
(8)~ProductionRunHealth: numerical-evidence claims must trace
to a successful \texttt{ProductionRun} graph node.

The output is a structured markdown report at
\texttt{papers/AutomatedReviews/<DATE>-internal-adversarial/<paper>.md}
with finding severity in \{BLOCKER, REQUIRED, RECOMMENDED\}.
\texttt{ReviewFinding} graph nodes are auto-extracted by
\texttt{scripts/build\_graph.py:extract\_review\_finding\_nodes}
with \texttt{FLAGS} edges to the affected paper or bundle, and
BLOCKER findings flip the affected
\texttt{ReadinessGate} status to \texttt{blocked}. Author-side fix
appends to a supersession ledger
(\texttt{docs/review\_finding\_supersessions.json}); on next graph
build the supersession overrides the finding's status to
\texttt{fixed} or \texttt{accepted}; \texttt{validate.py --check
readiness\_submission\_gate} re-evaluates the bundle.

The bundle-level extension of the pattern was added in Phase~6i
Wave~7. The orchestrator
\texttt{scripts/review\_runner.py --bundle <X> --prep-brief}
emits a per-bundle review brief; the agent walks every source
paper assigned to the bundle (per
\texttt{docs/PAPER\_DRAFT\_MAPPING.md}) plus the bundle's own
draft; output goes to
\texttt{papers/AutomatedReviews/<DATE>-bundle-stage13/<X>.md}.
The per-bundle anchor list at
\texttt{docs/agents/claims-reviewer-bundle-prompts.md} carries the
load-bearing claims and citations the bundle's review must verify.
Cross-bundle cluster mismatches (\texttt{cross\_bundle: true}
clusters in \texttt{cluster\_bundle\_index.json}) emit
\texttt{BundleConsistencyMismatch} findings.

The review-state semantics were hardened in June~2026 with a
\emph{recorded-review} requirement. A bundle's readiness verdict
can reach green only when three conditions hold simultaneously:
zero open blocker-class findings, at most five open advisories,
\emph{and} a recorded fresh-context Stage-13 review (the review
event itself, with its date and report path, captured in the
bundle's \texttt{bundle\_metadata.json}). The third condition is
the substantive one. Before the change, a bundle with no recorded
findings was indistinguishable from a bundle that had been
reviewed and passed; ``no findings recorded'' is not evidence of
review, it is evidence of nothing. The hardened semantics
introduce an explicit \emph{yellow (unreviewed)} verdict class for
bundles with zero open blockers but no recorded fresh-context
review, so the readiness heatmap cannot silently launder the
absence of review into a pass.

The June~2026 external-review remediation also surfaced the most
instructive citation-integrity lesson of the program to date. The
pipeline's citation machinery to this point (the primary-source
cache, Invariant~11, registry identifier resolution, and the
bibliographic-title drift detector) validates citation
\emph{metadata}. The remediation identified a failure class that
metadata validation cannot catch by construction: a citation can
be metadata-valid (the identifier resolves, the primary source
is cached, the bibliographic data are correct) while the
\emph{claim the prose attributes to it} is fabricated. Three
instances were found in roughly ten spot checks of load-bearing
attributions. In one, a draft attributed a quantitative
susceptibility-range estimate to a cached source that is a
unitarity proof in a lattice-gravity framework and contains no
such analysis. In another, the prose attributed an observational
exclusion to the wrong observables, at the wrong significance,
with the wrong year, relative to the cached paper. In the third,
the citation registry's own change history carried a fabricated
title and wrong journal coordinates for a real paper whose actual
framing is conditional rather than a no-go. All three passed every
metadata-level check, by construction: the sources are real,
cached, and correctly identified; the attributed \emph{content}
was wrong.

The deterministic response shipped in the same remediation as
three new validation checks. The first
(\texttt{validate.py --check axiom\_count\_prose\_consistency})
scans every paper draft for axiom-count claims and fails any
present-tense claim that disagrees with the live machine-generated
count, motivated by drafts still asserting the single-axiom
posture after the last project-local axiom was retired into a
tracked-hypothesis \texttt{Prop} on 2026-05-19; the calibration
run alone surfaced sixteen stale prose sites across eleven drafts.
The second
(\texttt{validate.py --check prose\_theorem\_reference\_coverage})
requires every Lean-identifier-shaped reference in monospaced
bundle prose to resolve to a live declaration in the Lean
dependency export, with a documented per-instance waiver mechanism
(this paper's deliberately historical reference to the disproved
folklore bound of Section~\ref{sec:case2} is the canonical
waiver). The third
(\texttt{validate.py --check theorem\_name\_embedded\_citations})
covers declaration names that embed author--year tokens: when a
paper's prose cites such a declaration, the citing paper's
bibliography must contain a matching entry, so a reader can
actually follow the citation the name encodes. The same
remediation also replaced a manually maintained deep-paper
headline count of cross-platform Lean-theorem reuse with a
deterministic scripted count
(\texttt{scripts/count\_theorem\_reuse.py}); the manual figure of
$109/119$ corrected to $106/114 \approx 93\%$, with both the
numerator and the denominator of the manual count having drifted.

Deterministic checks, however, guard \emph{presence and
consistency}; only reading the source against the claim verifies
\emph{support}. The systematic response is a content-level
attribution sweep (repository document
\texttt{docs/roadmaps/AttributionContentSweep\_Roadmap.md}). For
every load-bearing citation--claim pair on a bundle's anchor list,
a fresh-context agent reads the cached primary source and issues a
dated verdict: supported; supported-with-caveat (the paper must
then carry the caveat); not-supported (prose and registry
corrected, with a ledger-style record); or source-missing (fetch,
then verify). Each bundle visit pairs the content pass with a
fresh-context Stage-13 re-invocation, works falsification-critical
bundles first and the flagship last (so the flagship inherits its
siblings' corrections), and installs the register-derived
disclosure block of Section~\ref{sec:provenance}. The sweep's exit
criterion is every bundle green under the recorded-review
semantics above. The base rate that motivated the sweep
(roughly three fabricated attributions per ten spot-checked) is
exactly the kind of number a verification program must be willing
to publish about itself: the metadata layer the program had been
treating as citation integrity is necessary but not sufficient.

%% ─────────────────────────────────────────────────────────────────
%% §11 Phase 6f W1-W6 algebraic-GR substrate roster
%% ─────────────────────────────────────────────────────────────────

\section{Algebraic-General-Relativity Substrate Roster}
\label{sec:phase6f_substrate}

% TODO: substantive draft. Substrate to draw from:
%  - lean/SKEFTHawking/{ADMFormalism, AlgebraicRiemann,
%    SchwarzschildKerrSchild, DeSitterAdSLambdaVacuum,
%    TetradMetricEquivalence, ...}.lean
%  - Phase 6f Wave 1-6 memory files
%  - The algebraic-GR substrate roster:
%    + algebraic Riemann (Wave 1)
%    + Schwarzschild Kerr-Schild (Wave 2)
%    + de Sitter / AdS Lambda-vacuum biconditional (Wave 3)
%    + ADM Hamiltonian/momentum-constraint biconditionals (Wave 4)
%    + tetrad metric equivalence (Wave 5)
%    + Wave 6: tetrad formalism cross-bridges
%
% Section thesis: each of the six waves shipped a substrate-level
% formalization at the algebraic-general-relativity layer. These
% are infrastructure on which subsequent Phase 6 waves (BCH four
% laws, area theorem, Hawking-Penrose singularity theorem) build.
% The bundle paper presents this as a sidebar to the main
% methodology narrative.

Although this paper's primary contribution is methodological, the
verification pipeline has produced as a side product a body of
classical-general-relativity formalization at the algebraic
substrate level. We sketch the body here as a sidebar; full
treatment is the subject of a companion algebraic-general-relativity
backbone paper.

Phase~6f Waves~1 through~6 produced \phaseSixfWonesixTotal{}~theorems
(all substantive in the current build; the
substantive--vs--total distinction is the
\texttt{PLACEHOLDER\_THEOREMS} registry per Invariant~9 of
Section~\ref{sec:pipeline}) across six modules at the
algebraic (not differential-geometric)
level of GR. The algebraic level is the appropriate working language
for substrate-emergent gravity, where the metric, Christoffel
symbols, and curvature tensors are defined component-wise on
$\mathrm{Fin}\,4 \to \mathrm{Fin}\,4 \to \cdots \to \mathbb{R}$
arrays rather than via fiber-bundle abstractions; downstream
black-hole-thermodynamics work consumes the algebraic theorems as
Props rather than re-deriving them on a manifold. Six waves shipped
in sequence:

\begin{itemize}
\item \textbf{Wave~1 (\texttt{Curvature.lean}, \curvatureThms{}~theorems):}
algebraic Riemann tensor $R^\rho_{\sigma\mu\nu}$, the
\texttt{FirstBianchi} predicate, antisymmetry of the lower index pair,
and the pair-symmetry corollary $R_{\rho\sigma\mu\nu} =
R_{\mu\nu\rho\sigma}$.

\item \textbf{Wave~2 (\texttt{EinsteinTensor.lean}, \einsteinTensorThms{}~theorems):}
algebraic Einstein tensor $G_{\mu\nu} := \mathrm{Ric}_{\mu\nu} -
\tfrac{1}{2}R\, g_{\mu\nu}$, the dimension-4 trace identity
$G^\mu_\mu = -R$, the vacuum-characterization biconditional
$G_{\mu\nu} = 0 \iff \mathrm{Ric}_{\mu\nu} = 0$, and the de Sitter
$\Lambda$-vacuum cross-bridge $G + \Lambda g = 0 \iff \Lambda = 3K$
(constant-$K$ specialization).

\item \textbf{Wave~3 (\texttt{EnergyConditions.lean}, \energyConditionsThms{}~theorems):} pointwise energy-condition predicates (NEC, WEC, SEC,
DEC) on the algebraic stress-energy tensor and Ricci-curvature
specializations.

\item \textbf{Wave~4 (\texttt{ExactSolutions.lean}, \exactSolutionsThms{}~theorems):} four canonical exact solutions as algebraic consumers
of the Wave~1--3 substrate. Minkowski ($K=0$, $\Lambda=0$); de
Sitter ($K>0$, Hubble$^2 = K$); anti-de Sitter ($K<0$, AdS radius
$\ell = 1/\sqrt{-K}$); Schwarzschild in Kerr-Schild form
$g = \eta + (2M/r)\,l\otimes l$, with explicit cross-bridges to
Wave~1's Riemann tensor and downstream black-hole-thermodynamics
hooks.

\item \textbf{Wave~5 (\texttt{ADMFormalism.lean}, \admFormalismThms{}~theorems):}
ADM 3+1 decomposition. Hamiltonian constraint
$\square g = 2(R_{(3)} + K_{ij}K^{ij} - K^2)$, momentum constraint,
extrinsic curvature trace, with Minkowski and de Sitter explicit
witnesses.

\item \textbf{Wave~6 (\texttt{TetradFormalism.lean}, \tetradFormalismThms{}~theorems):} tetrad-induced metric symmetry, Minkowski tetrad
witness, determinant identity, and cross-bridges to the
Einstein-Cartan extension module (Phase~6e).
\end{itemize}

Total Phase~6f W1-W6 yield: \phaseSixfWonesixTotal{}~theorems,
all substantive (zero placeholder markers in the current build),
zero sorry, zero new axioms. The Phase~6f audit summary
is preserved in the program's per-phase literature cache,
maintained alongside the repository; downstream papers
(Bardeen-Carter-Hawking four laws, area theorem,
Hawking--Penrose singularity theorem) reference these modules as
load-bearing imports rather than as documentation. The
methodology point: a verification pipeline producing a sustained
research program inevitably accretes infrastructure that is itself
publishable; the substrate paper is structurally derivative of the
methodology, not parallel to it.

%% ─────────────────────────────────────────────────────────────────
%% §12 Phase 6f W7 Riemannian connection sidebar
%% ─────────────────────────────────────────────────────────────────

\section{The Riemannian Connection Sidebar}
\label{sec:riemannian_sidebar}

% TODO: substantive draft. Substrate to draw from:
%  - papers/paper44_riemannian_connection/paper_draft.tex
%  - lean/SKEFTHawking/{LeviCivita, BundleRiemann,
%    BundleRiemannAux, ...}.lean
%  - The session-by-session expansion (W7 sessions 1-3 SHIPPED;
%    sessions 4-5 in flight at memory write time)
%  - The mkHom₃ + bundled CovariantDerivative.riemann construction
%  - Bundle-level first/second Bianchi via Jacobi
%  - The Mathlib upstream PR target
%
% Section thesis: the Riemannian connection sidebar documents the
% structural completion of the curvature gap above the Bonn API.
% Wave 7 Sessions 1-3 shipped the Lorentzian metric typeclass +
% Levi-Civita Christoffel uniqueness + full coordinate Riemann +
% algebraic first Bianchi + differential second Bianchi machinery.
% Wave 8 Sessions 1-3 SHIPPED add the bundle-level Riemann
% curvature consumer of Bonn IsCovariantDerivativeOn + mlieBracket
% API. The sidebar is sized for I1's methodology paper, not for
% the full 100+ pp algebraic-GR backbone paper.

The algebraic substrate of Section~\ref{sec:phase6f_substrate}
operates pointwise on real-valued tensor components. A neighboring
question (equally classical, equally ``in the textbooks'')
is whether the connection-level structure on a smooth manifold
admits a parallel formalization. The Phase~6f Wave~7 catch-up
shipped the connection-level layer; Wave~8 lifted it from
component-coordinate scope to Bonn's bundle-level
\texttt{IsCovariantDerivativeOn} API. We sketch the result here,
keeping detail proportionate to a methodology-paper sidebar; the
a companion paper currently in preparation as a project-internal
draft treats the construction as its primary subject.

The starting point is the indefinite-signature metric typeclass
\texttt{IsLorentzianMetric}. Mathlib's existing infrastructure on
Riemannian-bundle metrics
(\texttt{Mathlib.Geometry.Manifold.IsRiemannianBundle}, due to
S.~Gou\"ezel) is positive-definite by design; Wave~7's contribution
is the indefinite-signature analogue.
Our typeclass is a field-level predicate (signature
$(-, +, +, +)$), with Minkowski as a primitive witness
(\texttt{$\eta$\_isLorentzianMetric}). Wave~7 (modules
\texttt{LorentzianMetric.lean} + \texttt{RiemannianConnection.lean},
\phaseSixfWsevenTotal{}~theorems) shipped the Christoffel-quadratic
Riemann tensor antisymmetry (in the lower index pair, separate
from the Wave~1 algebraic version which constructs Riemann from
the abstract type rather than from the metric), the
Koszul-formula closed form for the Levi-Civita Christoffel
symbols, and a uniqueness predicate at the
Levi-Civita-Christoffel level. A predicate-level falsifier
asserts that a Riemannian-signature metric cannot satisfy
\texttt{IsLorentzianMetric}.

Wave~8 Sessions~1 through~5 (modules
\texttt{RiemannCoordinate.lean},
\texttt{RiemannDifferentialBianchi.lean},
\texttt{BundleRiemannAux.lean}, \texttt{BundleRiemann.lean},
\texttt{LeviCivita.lean}; \phaseSixfWeightTotal{}~theorems total:
\riemannCoordinateThms{} + \riemannDifferentialBianchiThms{} +
\bundleRiemannAuxThms{} + \bundleRiemannThms{} + \leviCivitaThms{} by
module) shipped the coordinate-form Riemann tensor in linear
plus quadratic decomposition, the algebraic first Bianchi
identity, the differential second Bianchi identity, the
bundle-level Riemann tensor consuming Bonn's
\texttt{IsCovariantDerivativeOn} + \texttt{mlieBracket} API, the
cyclic Jacobi identity for the manifold Lie bracket (Mathlib
ships only the Leibniz form), the bundle-level first Bianchi
identity under torsion-freeness, and the substantive bundle-level
Levi-Civita uniqueness theorem via the Koszul-bilinear-form
derivation $2g(\nabla_X Y, Z) = X g(Y,Z) + Y g(Z,X) - Z g(X,Y) +
g([X,Y], Z) - g([X,Z], Y) - g([Y,Z], X)$.

The bundle-level Levi-Civita uniqueness theorem is the
load-bearing result of Wave~8: it lifts the Wald \S3.1 /
Kobayashi-Nomizu IV.2.2 textbook construction to a
proof-assistant-checked form, using \texttt{FiberBundle.extend}
plus bilinearity non-degeneracy. The Phase~6f W7--W8 substrate
collectively addresses two Mathlib gaps: the
indefinite-signature metric typeclass, and Riemann/Ricci/scalar
curvature on a connection (out-of-scope for the
Massot--Rothgang--Macbeth Half-ERC proposal explicitly per their
publicly visible roadmap). Two upstream Mathlib pull requests
are in preparation; the current pin is a vendored copy in
\texttt{lean/SKEFTHawking/}.

The sidebar's methodology point is that the verification pipeline
produces \emph{textbook formalization} as a byproduct of doing
applied physics in Lean. The five-page companion paper that
documents this work (in preparation as a project-internal draft;
target journal a Mathlib-aligned formal-verification venue) is a
publication target in its own right, unrelated to the SK-EFT
Hawking program except in that the Levi-Civita uniqueness
statement was needed by the Bardeen-Carter-Hawking four-laws
result of paper~L3 / D3 §8.

%% ─────────────────────────────────────────────────────────────────
%% §13 Phase 6g W5 structural-Prop scoping pattern
%% ─────────────────────────────────────────────────────────────────

\section{Structural-Prop Scoping Under Mathlib-Dependency Fallback}
\label{sec:structural_prop_scoping}

% TODO: substantive draft. Substrate to draw from:
%  - lean/SKEFTHawking/{HawkingPenroseSingularityCurveTheoretic,
%    AreaTheoremCurveTheoretic, RiccatiComparison, ...}.lean
%  - Phase 6g Wave 5 memory file
%  - The 1-D wave equation Fourès-Bruhat distillation case
%  - The pattern: when Mathlib lacks the PDE infrastructure
%    (Cauchy-Kowalevski, energy estimates), structure the
%    substantive content as a tracked Prop bundle that the
%    project-local proof discharges modulo the missing infra,
%    and document the missing infra as a "Mathlib upstream
%    target" rather than as a sorry
%
% Section thesis: when a substantive Lean theorem cannot be
% closed because the upstream PDE / measure-theoretic /
% functional-analytic machinery is unavailable in Mathlib, the
% structurally-correct response is not to invent the missing
% theory in-project. It is to scope the substantive content as a
% tracked Prop bundle, prove every other hypothesis in the
% bundle, and let the missing-infra Prop be the gating hypothesis
% that downstream proofs consume. This avoids in-project
% reinvention of Mathlib while keeping the substantive content
% Lean-checkable.

A recurring failure mode in the project's first two years was
in-project reinvention of upstream mathematical machinery whenever a
target theorem depended on infrastructure absent from Mathlib. The
Cauchy problem for the Einstein vacuum equations, the
Cauchy--Kowalevski theorem, energy estimates for hyperbolic systems,
and distributional Sobolev embedding all entered our Lean wishlist at
various points. Each attempt to formalize from scratch consumed
two-to-eight weeks of effort and produced incomplete substrate that
the next wave was reluctant to reuse. The cumulative cost of these
attempts ran into the months.

The Phase~6g Wave~5 finding, refined at Waves~9 and~10, was that
\emph{the structurally-correct response is not in-project
reinvention}. The pattern is straightforward: when a theorem
depends on Mathlib-absent infrastructure, scope the missing
infrastructure as a tracked-\texttt{Prop} bundle, prove every other
hypothesis in the bundle within project scope, and let the
missing-infrastructure \texttt{Prop} be the gating hypothesis that
downstream proofs consume as an assumption. The result is a
substantive theorem of the form ``modulo Mathlib provides
\texttt{IsLocallyWellPosed}, the Cauchy problem is well-posed and
satisfies $X$.''

Wave~5 itself
(\texttt{lean/SKEFTHawking/CauchyProblem.lean}) shipped the
Einstein-vacuum Cauchy problem in this form. Wave~9 produced four
substantive curve-theoretic modules
(\texttt{RiccatiComparison.lean}, \texttt{MazurSigmaModelRigidity.lean},
\texttt{HawkingPenroseSingularityCurveTheoretic.lean},
\texttt{AreaTheoremCurveTheoretic.lean}) yielding \phaseSixgWnineTotal{}~theorems
(\riccatiComparisonThms{} + \mazurSigmaModelRigidityThms{} +
\hawkingPenroseCurveThms{} + \areaTheoremCurveThms{} by module) on
Raychaudhuri-Riccati comparison, the
$\sigma$-model 1D rigidity (Mazur's lemma:
$f' \geq 0 \wedge f(0) = f(L) \implies f \equiv f(0)$), the
curve-theoretic Hawking--Penrose singularity theorem with a
de~Sitter strong-energy-condition counterexample, and area-evolution
ODE monotone rigidity. Wave~10 closed the Phase~6g cycle with the
distillation case study: the one-dimensional wave equation
$\partial^2 u / \partial t^2 = c^2\, \partial^2 u / \partial x^2$
reduced from the Einstein vacuum equations in harmonic gauge.
\texttt{WaveEquation1D.lean} ships an
\texttt{IsClassicalWaveSolution} predicate, two explicit witnesses
(constant solution; linear traveling-wave solution),
the linearity corollary
(\texttt{add\_isClassicalWaveSolution}: sum of solutions is a
solution), and the cross-bridge to the structural-\texttt{Prop}
\texttt{IsLocallyWellPosed} that downstream Cauchy-problem theorems
consume.

The discipline is what makes substantive Lean content tractable in
a research domain whose mathematical foundation is not yet in
Mathlib. The cost of admitting a structural assumption with a clear
provenance link to an upstream target is small; the cost of
reinventing the upstream target is large; the cost of \emph{mixing}
in-project partial reinvention with downstream theorems that
consume it is unbounded, because partial reinvention rots
asymmetrically when Mathlib eventually catches up. The pattern's
emergence in Phase~6g, midway through the project, is itself a
methodology lesson: structural-\texttt{Prop} scoping is the right
default from the first wave; we encoded the discipline only after
two months of recoverable cost.

%% ─────────────────────────────────────────────────────────────────
%% §14 Lessons + future work
%% ─────────────────────────────────────────────────────────────────

\section{Lessons and Future Work}
\label{sec:lessons}

We close with four lessons drawn from a year of running the
pipeline at the scale described.

\textbf{Continuous verification is cheaper than retrofitted
verification, by a margin large enough to be a workflow choice,
not a methodology preference.} Each of the four failure modes we
opened with (a numerical claim that traces to a non-existent
code path, a ``formally verified'' claim whose theorem is a
placeholder, a citation that resolves to the wrong paper, a
parameter that has drifted from its primary source) would
have shipped to draft state under the pre-pipeline workflow. Each
was caught, in this project, by the pipeline's continuous
checks before it could be cited downstream. The cost of running
the checks is on the order of seconds per stage; the cost of
finding and fixing a propagated error after publication is on the
order of weeks per finding, in the best case.

\textbf{Automated theorem proving is a force multiplier on
formalization work, but only when the formalization itself is
well-decomposed.} A monolithic fifty-term proof goal is
intractable to both manual proof and to Aristotle. A decomposed
nine-deep structure with each sublemma at twelve terms or fewer is
tractable to both, and Aristotle closes the sublemmas in a single
priority batch. The decomposition discipline is the pipeline's
load-bearing investment: it pays back in throughput at every
subsequent wave.

\textbf{Adversarial review is a separate role from authorship,
and the separation is load-bearing.} A reviewer in the same
context as the author has, by construction, internalized the
author's prior justifications. The reviewer's findings will, by
construction, be limited to what the author's context permits.
The fresh-context adversarial-reviewer pattern (the reviewer
agent runs in a context window that has not seen the author's
work) is the structural fix. The author fixes; the author
re-invokes; iteration converges to clean.

\textbf{Metadata-valid is not content-true.} The citation
machinery of Invariant~11 (cached primary sources, resolving
identifiers, title-drift detection) structurally prevents
hallucinated \emph{references}. It does not prevent fabricated
\emph{attributions}: prose ascribing to a real, cached, correctly
cited source a claim the source does not make. At a measured base
rate of roughly three per ten spot-checked load-bearing
attributions (Section~\ref{sec:adversarial}), this is the failure
mode that survived every deterministic layer of the pipeline. The
response (deterministic prose-consistency checks where the
property is mechanizable, fresh-context content reading where it
is not) is the pipeline's newest layer, and there is no reason
to believe it is the last.

Three pipeline extensions matured during the year under review
and are now load-bearing. First, a meta-process quality-improvement register
(Stage~14) that surfaces systemic patterns across waves; the seed
event is the April~2026 external-adversarial-review round that
exposed a thirteen-dimensional problem space that internal
validation could not detect by construction. Second, the
publication-bundle architecture that consolidates per-wave drafts
into the publication targets that frame enumerates while preserving sentence-level
provenance; this paper is drafted under that architecture. Third, an automated cross-bundle consistency check
that walks a cross-bundle cluster registry and flags member-
sentence drift across bundle boundaries, preventing the same fact
from being stated differently in two papers.

The pipeline is open-source, domain-agnostic, and immediately
applicable in any research domain where the connection between
computation and published claims must be trustworthy. We
specifically encourage adoption in any setting where the cost of a
correctness failure post-publication exceeds the cost of running
continuous verification: a condition that holds in
pharmaceutical software (FDA SaMD algorithmic vigilance),
verified-risk financial modeling, climate modeling, and any other
research domain where the audit trail from code to claim is itself
the deliverable.

%% ─────────────────────────────────────────────────────────────────
%% Code availability
%% ─────────────────────────────────────────────────────────────────

\section*{Code Availability}

The complete pipeline described in this paper (the Python
numerics, the Lean~4 formalization, the validation and
graph-construction scripts, \texttt{scripts/}, the canonical
process documents, \texttt{docs/}, the automated-review reports,
\texttt{papers/AutomatedReviews/}, and the source of this paper
itself) is maintained in a single public repository at
\url{https://github.com/NetRxn/SK_EFT_Hawking}. Monospaced file
paths throughout this paper are paths from the repository root;
templated path segments in angle brackets (e.g.,
\texttt{run\_<timestamp>}) denote per-instance artifacts.

%% ─────────────────────────────────────────────────────────────────

\begin{thebibliography}{99}
\bibitem{physlean} J.~Tooby-Smith, ``PhysLean: Lean~4 for physics,''
\url{https://github.com/HEPLean/PhysLean} (2024).
\bibitem{compcert} X.~Leroy, ``Formal verification of a realistic
compiler,'' \emph{Comm.\ ACM}\ \textbf{52}(7), 107--115 (2009).
\bibitem{sel4} G.~Klein et al., ``seL4: Formal verification of an OS
kernel,'' \emph{Proc.\ ACM SOSP} (2009).
\bibitem{leancipilot} P.~Song, K.~Yang, and A.~Anandkumar, ``Lean
Copilot: Large Language Models as Copilots for Theorem Proving in
Lean,'' arXiv:2404.12534 (2024).
\bibitem{alphaproof} Google DeepMind, ``AI achieves silver-medal
standard solving International Mathematical Olympiad problems,''
(2024).

\bibitem{CrossleyGloriosoLiu2017} M.~Crossley, P.~Glorioso, and
H.~Liu, ``Effective field theory of dissipative fluids,''
\emph{JHEP}\ \textbf{09} (2017) 095, arXiv:1511.03646.

\bibitem{GloriosoCrossleyLiu2017} P.~Glorioso, M.~Crossley, and
H.~Liu, ``Effective field theory for dissipative fluids (II):
classical limit, dynamical KMS symmetry and entropy current,''
\emph{JHEP}\ \textbf{09} (2017) 096, arXiv:1701.07817.

\bibitem{GloriosoLiu2018} P.~Glorioso and H.~Liu, ``The second law
of thermodynamics from symmetry and unitarity,'' arXiv:1612.07705
(2016).

\bibitem{GloriosoLiu2018TASI} P.~Glorioso and H.~Liu, ``Lectures
on non-equilibrium effective field theories and fluctuating
hydrodynamics,'' \emph{PoS}\ \textbf{TASI2017} (2018) 008,
arXiv:1805.09331.

\bibitem{JainKovtun2024} A.~Jain and P.~Kovtun, ``Schwinger--Keldysh
effective field theory for stable and causal relativistic
hydrodynamics,'' \emph{JHEP}\ \textbf{01} (2024) 162,
arXiv:2309.00511.
\bibitem{CsikorFodorHeitger1999} F.~Csikor, Z.~Fodor and
J.~Heitger, ``Endpoint of the hot electroweak phase transition,''
\emph{Phys.\ Rev.\ Lett.}\ \textbf{82} (1999) 21,
arXiv:hep-ph/9809291.

\bibitem{KLRS1996} K.~Kajantie, M.~Laine, K.~Rummukainen and
M.~Shaposhnikov, ``Is there a hot electroweak phase transition at
$m_H \gtrsim m_W$?,'' \emph{Phys.\ Rev.\ Lett.}\ \textbf{77} (1996)
2887, arXiv:hep-ph/9605288.

\end{thebibliography}


%% ─── Lifted from _phase6n_W2b_lean_only (§8) — 2026-05-06T15:42:06Z ───
%% \section{Quantum Crooks no-go (Perarnau-Llobet)}
%% \label{sec:i1-_phase6n_W2b_lean_only}

%% Insertion point hint from PAPER_DRAFT_MAPPING.md: §8
%% Lift notes: D.2 sidebar: Phase 6n W2b methodology case — typeclass-connections cross-bridge pattern; QuantumCrooks/SKEFTConnection IsReservoirCoupled non-vacuity strengthening
%% TODO: lift content from papers/_phase6n_W2b_lean_only/paper_draft.tex
%% TODO: ensure all numerical claims trace via formulas.py / counts.tex
%% TODO: ensure all citations have primary-source cache entries

%% (Section heading + label commented out 2026-05-06 Session 2: empty post-bibliography stub from bundle_append.py default-insertion; bookkeeping anchor preserved as LaTeX comment, no rendered content. Append-log entry retains source-contribution registration.)
%% ─── End lift from _phase6n_W2b_lean_only ───


%% ─── Lifted from _phase6o_W2a_lean_only (§11) — 2026-05-06T15:42:10Z ───
%% \section{APS $\eta$-invariant for analog horizons: substrate-discovery methodology sidebar (Phase 6o W2a)}
%% \label{sec:i1-_phase6o_W2a_lean_only}

%% Insertion point hint from PAPER_DRAFT_MAPPING.md: §11
%% Lift notes: D.2 sidebar: Phase 6o W2a methodology case — substrate-discovery pattern (predicate-level operationalization at Wave 2a.2 + per-substrate concrete computation at Wave 2a.3-5; ³He-A unique non-degenerate cell)
%% TODO: lift content from papers/_phase6o_W2a_lean_only/paper_draft.tex
%% TODO: ensure all numerical claims trace via formulas.py / counts.tex
%% TODO: ensure all citations have primary-source cache entries

%% (Section heading + label commented out 2026-05-06 Session 2: empty post-bibliography stub from bundle_append.py default-insertion; bookkeeping anchor preserved as LaTeX comment, no rendered content. Append-log entry retains source-contribution registration.)
%% ─── End lift from _phase6o_W2a_lean_only ───


%% ─── Lifted from _phase6o_W3a_lean_only (§7) — 2026-05-06T15:42:59Z ───
%% \section{ETH-$\alpha$ refutation-tableau methodology sidebar (Phase 6o W3a)}
%% \label{sec:i1-_phase6o_W3a_lean_only}

%% Insertion point hint from PAPER_DRAFT_MAPPING.md: §7
%% Lift notes: D.2 sidebar: Phase 6o W3a methodology case — parallel-axiomatization productive-value Aristotle-driven tableau pattern (5 candidate ETH axiomatizations + 3 concrete refutation theorems)
%% TODO: lift content from papers/_phase6o_W3a_lean_only/paper_draft.tex
%% TODO: ensure all numerical claims trace via formulas.py / counts.tex
%% TODO: ensure all citations have primary-source cache entries

%% (Section heading + label commented out 2026-05-06 Session 2: empty post-bibliography stub from bundle_append.py default-insertion; bookkeeping anchor preserved as LaTeX comment, no rendered content. Append-log entry retains source-contribution registration.)
%% ─── End lift from _phase6o_W3a_lean_only ───

\end{document}
